\BourbakiExerciseGroup

All rings under consideration are assumed to be commutative unless explicitly stated otherwise.

\begin{enumerate}
\def\labelenumi{\arabic{enumi})}
\tightlist
\item
  Let A be a totally ordered ring and B an additive subgroup of A which is closed under multiplication.
\end{enumerate}

\begin{enumerate}
\def\labelenumi{\alph{enumi})}
\item
  An element \(x \in A\) is said to be infinitely large with respect to B if \(|y| < |x|\) for all \(y \in B\) . Show that the set \(F(B)\) of elements of A which are not infinitely large with respect to B is a subring of A containing B.
\item
  An element \(x \in A\) is said to be infinitely small with respect to B if, for all y \textgreater{} 0 in B, we have \(|x| < y\) . If, for all y \textgreater{} 0 in B, there exists \(z \in B\) such that 0 \textless{} z \textless{} y, then the set of elements of A which are infinitely small with respect to B is a subpseudo-ring I(B) of A. If in addition, for every pair of elements y, z of B with 0 \textless{} y \textless{} z, there exists \(x \in B\) with 0 \textless{} xz \textless{} y, then I(B) is an ideal in the ring F(B).
\end{enumerate}

\begin{enumerate}
\def\labelenumi{\arabic{enumi})}
\setcounter{enumi}{1}
\item
  Let A be a totally ordered ring and n the set of nilpotent elements of A (which is an ideal of A). Every element of A not belonging to n is infinitely large with respect to n; the quotient ring A/n is totally ordered (VI, p.~30, Ex. 4) and is an integral domain.
\item
  In the field \(\mathbf{Q}\) , let \(\mathbf{P}\) be the set consisting of 0 and the rational numbers \(\geqslant 1\) (in the usual ordering). Show that \(\mathbf{P}\) satisfies axioms \((\mathrm{AP}_{\mathrm{I}})\) , \((\mathrm{AP}_{\mathrm{II}})\) and \((\mathrm{AP}_{\mathrm{III}})\) .
\item
  \begin{enumerate}
  \def\labelenumii{\alph{enumii})}
  \tightlist
  \item
    Let K be a field, and P a subset of K satisfying conditions \((\mathrm{AP}_{\mathrm{I}})\) , \((\mathrm{AP}_{\mathrm{II}})\) and \((\mathrm{AP}_{\mathrm{III}})\) , and such that \(K^{2} \subset P\) (where \(K^{2}\) denotes the set of squares of elements of K); show that if x \textgreater{} 0 in the ordering defined by P, then \(x^{-1} > 0\) ; deduce that the set \(K_{+}^{*}\) of elements \textgreater{} 0 in K is a subgroup of the multiplicative group of K.
  \end{enumerate}
\end{enumerate}

\begin{enumerate}
\def\labelenumi{\alph{enumi})}
\setcounter{enumi}{1}
\item
  In the field \(\mathbf{Q}\) , the only set \(\mathbf{P}\) which satisfies the conditions of \(a\) ) is the set of rational numbers \(\geqslant 0\) (in the usual ordering).
\item
  * Let P be a subset of K satisfying \((\mathrm{AP}_{\mathrm{I}})\) , \((\mathrm{AP}_{\mathrm{II}})\) and \((\mathrm{AP}_{\mathrm{III}})\) , such that \(1 \in P\) and such that x \textgreater{} 0 implies \(x^{-1} > 0\) in the ordering defined by P. Show that \(y^{2} \geqslant z^{2}\) implies \(y \geqslant z\) , for positive y and z. Deduce that, for arbitrary y \textgreater{} 0, we have
\end{enumerate}

\[
\frac {1}{2} (y + y ^ {- 1}) \geqslant 1 - n ^ {- 1}
\]

for every integer \(n > 0\) (notice that \((y + y^{-1})^{2m} \geqslant \binom{2m}{m}\) for every integer \(m > 0\) ). Deduce that, if the ordered additive group structure defined by P is archimedean (VI, p.~35, Ex. 31), then \((y - z)^2 \geqslant 0\) for every pair of elements \(y, z\) of P; if \(K'\) is the subfield of K generated by P then \(x^2 \geqslant 0\) for all \(x \in K'\) . *

\(d) * \text{Let } K = R(X) \text{ be the field of rational functions in one indeterminate over } R; \text{ let } P \text{ be the set consisting of 0 and the rational functions } u \in K \text{ such that } u(t) \text{ is defined and } >0 \text{ for every real number } t. \text{ Show that } P \text{ satisfies the conditions of } c) \text{ and generates } K, \text{ but that there exist elements } v \in K \text{ such that } v^2 \notin P. *\)

\begin{enumerate}
\def\labelenumi{\arabic{enumi})}
\setcounter{enumi}{4}
\tightlist
\item
  Let A be a lattice ordered ring. A convex ideal (VI, p.~30, Ex. 4) of A is called irreducible if it is not the intersection of two convex ideals distinct from itself.
\end{enumerate}

\begin{enumerate}
\def\labelenumi{\alph{enumi})}
\tightlist
\item
  Show that the intersection of the irreducible convex ideals of A is 0 (if \(a \in \mathbf{A}\) is nonzero, consider a convex ideal of A maximal among those which do not contain \(a\) ). Deduce that A is isomorphic to a subring \(\mathbf{A}'\) of a product \(\prod_{i} \mathbf{A}_i\) , such that each \(\mathbf{A}_i\) is lattice ordered,
\end{enumerate}

\(\mathrm{pr}_{\iota}(\mathbf{A}^{\prime})=\mathbf{A}_{\iota}\) for all \(\iota\) and the ideal \(\{0\}\) is irreducible in \(\mathbf{A}_{\iota}\) (as a convex ideal) for all \(\iota\) . b) Show that the following conditions are equivalent in A:

α) \(|xy| = |x| \cdot |y|\) for all x, y;

\(\beta) x \cdot \sup(y, z) = \sup(xy, xz)\) for all \(x \geqslant 0\) , \(y, z\) ;

\(\gamma) x \cdot \inf(y, z) = \inf(xy, xz)\) for all \(x \geqslant 0\) , \(y, z\) ;

\(\delta)\inf (y,z) = 0\) implies \(\inf (xy,z) = 0\) for all \(x\geqslant 0\) .

(To see that \(\gamma\) ) implies \(\delta\) ), observe that \(xy \leqslant y \cdot \sup (x, 1)\) ).

When these conditions are satisfied, then A is said to be strongly lattice ordered.

\begin{enumerate}
\def\labelenumi{\alph{enumi})}
\setcounter{enumi}{2}
\tightlist
\item
  Show that a ring \(\mathbf{A}\) is strongly lattice ordered if and only if it is isomorphic to a subring of a product \(\prod_{i} \mathbf{A}_{i}\) of totally ordered rings. (Using \(a\) ), reduce to showing that if
\end{enumerate}

\{0\} is irreducible in a strongly lattice ordered ring A, then A is totally ordered; for this, note that if B is a strongly lattice ordered ring then for all \(b \in B\) the set \(m\) of \(x \in B\) such that \(\inf(|x|, |b|) = 0\) and the set \(n\) of \(y \in B\) such that \(|y| \leqslant |zb|\) for some \(z \in B\) are two convex ideals of B such that \(m \cap n = \{0\}\) .

\(d\) ) Show that in a strongly lattice ordered ring \(\mathbf{A}\) the relation \(\inf (x,y) = 0\) implies \(xy = 0\) ; for all \(z\in \mathbf{A}\) we have \(z^2\geqslant 0\) ; for all \(x,y\in \mathbf{A}\) we have \(xy\leqslant \sup (x^2,y^2)\)

\begin{enumerate}
\def\labelenumi{\alph{enumi})}
\setcounter{enumi}{4}
\tightlist
\item
  Let A be a lattice ordered ring with no nilpotent element \textgreater0; show that if \(\inf(x, y) = 0\) implies xy = 0 then A is strongly lattice ordered (check condition \(\delta\) ) of b)).
\end{enumerate}

\begin{enumerate}
\def\labelenumi{\arabic{enumi})}
\setcounter{enumi}{5}
\tightlist
\item
  \begin{enumerate}
  \def\labelenumii{\alph{enumii})}
  \tightlist
  \item
    Let K be a lattice ordered field. Show that the following conditions are equivalent:
  \end{enumerate}
\end{enumerate}

\(\alpha) x^{2} \geqslant 0\) for all \(x \in \mathbf{K}\) ;

\(\beta) x > 0\) implies \(x^{-1} > 0\) ;

γ) K is strongly lattice ordered (Ex. 5);

δ) K is totally ordered.

(To see that \(\beta\) ) implies \(\underline{\delta}\) ), notice that \(\beta\) implies \(xy \leqslant x^2 + y^2\) for \(x, y \in K\) ).

\begin{enumerate}
\def\labelenumi{\alph{enumi})}
\setcounter{enumi}{1}
\item
  Let K be the field \(\mathbf{Q}(\sqrt{2})\) obtained by adjoining \(\sqrt{2}\) to the field of rational numbers; as an additive group K is identified with \(Q \times Q\) via the bijection \(x + y \sqrt{2} \mapsto (x, y)\) ; show that if the product ordering on \(Q \times Q\) (where Q has the usual ordering) is carried over to K via this bijection, then K is lattice ordered but not strongly lattice ordered.
\item
  The field \(\mathbf{K} = \mathbf{Q}(\mathbf{X})\) of rational functions in one indeterminate X over Q is generated by its set of squares \(K^{2}\) . Show that the set P of sums of squares of elements of K defines a non-lattice ordering on K compatible with its ring structure, and such that u \textgreater{} 0 implies \(u^{-1} > 0\) .
\end{enumerate}

\begin{enumerate}
\def\labelenumi{\arabic{enumi})}
\setcounter{enumi}{6}
\item
  Every element of a commutative field K of characteristic \(\neq 2\) is a sum of squares if and only if -1 is a sum of squares (notice that every element of K is a difference of two squares).
\item
  Let A be a nonempty subset of a commutative field K of characteristic \(\neq 2\) ; then K is said to be A-orderable if no element of A is a sum of squares in K. We say that K is orderable if there exists a total ordering of K which is compatible with its ring structure.
\end{enumerate}

\begin{enumerate}
\def\labelenumi{\alph{enumi})}
\item
  Show that if K is A-orderable then it is orderable (Ex. 7), and hence of characteristic 0.
\item
  Show that pure extensions and algebraic extensions of odd degree of an A-orderable field K are A-orderable (argue as in Prop. 3 of VI, p.~23).
\item
  Let K be an A-orderable field and let b be an element of K not of the form ca - d, where \(a \in A\) and c and d are sums of squares in K. Show that the field \(\mathrm{K}(\sqrt{b})\) is A-orderable.
\item
  An A-orderable field K is called maximal if no proper algebraic extension of K is A-orderable. Show that a maximal A-orderable field K has the following properties: 1. K is pythagorean, in other words every sum of squares is a square; 2. no element of A is a square; 3. every element of K which is not a square has the form \(c^2 a - d^2\) , where \(a \in \mathbf{A}\) ; 4. every odd degree polynomial in K{[}X{]} has at least one root in K (use b) and c)). e) Show that if K satisfies the conditions of d) then every element algebraic over K has degree \(2^q\) over K for some q (argue as in Prop. 8 of VI, p.~26, by induction on the exponent of 2 in the degree of the element under consideration). Show that on the other hand no quadratic extension of K is A-orderable. Deduce that K is a maximal A-orderable field (use Galois theory and Prop. 12 of I, p.~77).
\item
  Let K be an A-orderable field and \(\Omega\) an algebraically closed extension of K. Show that there exists a maximal A-orderable field contained in \(\Omega\) and containing K.
\end{enumerate}

\begin{enumerate}
\def\labelenumi{\arabic{enumi})}
\setcounter{enumi}{8}
\tightlist
\item
  Let \(q > 0\) be a non-square natural number; let \(A\) denote the set \(\{-1, \sqrt{q}\}\) in the field \(K = Q(\sqrt{q})\) ; show that \(K\) is \(A\) -orderable. Show that there exists an extension \(E\) of \(K\) such that \(E\) is pythagorean, orderable, and that every odd degree polynomial over \(E\) has a root in \(E\) , but that \(E\) does not admit the structure of a maximal ordered field (consider a maximal \(A\) -orderable extension of \(K\) ).
\end{enumerate}

¶ 10) a) Let K be an orderable field and E a Galois extension of K. Show that either E is orderable or there exists an orderable algebraic extension F of K, contained in E, such that E is a quadratic extension of F (use Th. 6 of V, p.~70).

\begin{enumerate}
\def\labelenumi{\alph{enumi})}
\setcounter{enumi}{1}
\tightlist
\item
  Show that the polynomial \(X^{4} + 2\) is irreducible over Q, and that the extension K of Q of degree 4, obtained by adjoining a root of this polynomial to Q, contains no orderable subfield other than Q (find all the subfields of K by Galois theory).
\end{enumerate}

\begin{enumerate}
\def\labelenumi{\arabic{enumi})}
\setcounter{enumi}{10}
\tightlist
\item
  Let K be an ordered field and G a subfield of K. a) Show that \(x \neq 0\) is infinitely large with respect to G if and only if \(x^{-1}\) is infinitely small with respect to G. We say that K is comparable to G if there is no element of K which is infinitely large with respect to G (nor, consequently, any nonzero element of K which is infinitely small with respect to G). A necessary and sufficient condition for K to be comparable to its prime subfield Q is that K be archimedean (VI, p.~35, Ex. 31), * and hence that K be isomorphic to a subfield of R (VI, p.~36, Ex. 33c)).
\end{enumerate}

\begin{enumerate}
\def\labelenumi{\alph{enumi})}
\setcounter{enumi}{1}
\item
  Show that, in the ring \(F(G)\) of elements of K which are not infinitely large with respect to G, the set \(I(G)\) of elements infinitely small with respect to G is a maximal ideal; moreover the ordering on the quotient field \(K(G) = F(G)/I(G)\) induced from that of \(F(G)\) is compatible with the ring structure and is total.
\item
  Show that a class modulo \(I(G)\) can contain at most one element of G; deduce that the natural map from G to \(K(G)\) is an isomorphism from the ordered field G onto a subfield \(G'\) of \(K(G)\) , and that \(K(G)\) is comparable to \(G'\) .
\end{enumerate}

\begin{enumerate}
\def\labelenumi{\arabic{enumi})}
\setcounter{enumi}{11}
\item
  Let K be a maximal ordered field, let f be a polynomial over K and let a and b be two roots of f such that a \textless{} b and such that f has no roots between a and b. Show that if g is a rational function over K whose denominator does not vanish for \(a \leqslant x \leqslant b\) then the equation \(f(x)g(x) + f'(x) = 0\) has an odd number of solutions in \([a, b]\) (where each solution is counted with its multiplicity; use Prop. 5 of VI, p.~25). Deduce that if h is a rational function over K having a and b as roots and whose denominator does not vanish in \([a, b]\) , then the equation \(h'(x) = 0\) has at least one root in \([a, b]\) .
\item
  Let K be a maximal ordered field, let h be a rational function over K and let \([a, b]\) be a closed interval in which h is defined. Show that there exists \(c \in ]a, b[\) such that \(h(b) - h(a) = (b - a)h'(c)\) (use Ex. 12). Deduce that h is an increasing function in \([a, b]\) if and only if \(h'(x) \geqslant 0\) in this interval (to see that this condition is necessary, split the interval by the roots of \(h'(x) = 0\) ).
\item
  \begin{enumerate}
  \def\labelenumii{\alph{enumii})}
  \tightlist
  \item
    Let K be a maximal ordered field, let E be a subfield of K and let f be a polynomial over E; show that all the roots of f in K belong to F(E) (Ex. 11, b)) (use Prop. 4 of VI, p.~24). Deduce that if G is a subfield of K and E ⊂ K is an extension comparable to G (Ex. 11) then the set of elements of K algebraic over E is a maximal ordered field comparable to G.
  \end{enumerate}
\end{enumerate}

\begin{enumerate}
\def\labelenumi{\alph{enumi})}
\setcounter{enumi}{1}
\item
  Deduce from a) that, under the hypotheses of a), the field K(G) (Ex. 11, b)) is a maximal ordered field.
\item
  Let \(f\) be a polynomial of degree \(\geqslant 1\) over \(G\) . Show that \(f(t)\) is infinitely large with respect to \(G\) if and only if \(t\) is infinitely large with respect to \(G\) ; the element \(f(t)\) is infinitely small with respect to \(G\) if and only if \(t\) is congruent modulo \(I(G)\) to a root of \(f\) in \(K\) (split \(K\) into intervals using the roots of \(f\) and \(f'\) in \(K\) , and apply Ex. 13; notice that if \(x < t\) and if \(x \in K\) is not congruent modulo \(I(G)\) to \(t\) , then there exists \(y \in K\) such that \(x < y < t\) and \(y - x \in G\) ).
\end{enumerate}

\begin{enumerate}
\def\labelenumi{\arabic{enumi})}
\setcounter{enumi}{14}
\tightlist
\item
  Let E be an ordered field, let K be a subfield of E such that E is algebraic over K, and let R be a maximal ordered extension of K.
\end{enumerate}

\begin{enumerate}
\def\labelenumi{\alph{enumi})}
\item
  Among the elements of E - K, let \(x_0\) be one whose minimal polynomial \(f\) over K has the least possible degree. Show that \(f'(T)\) is a product of factors in R{[}T{]} of the form \((T - a_i)^2 + c_i^2\) ( \(a_i, c_i \in \mathbb{R}\) ), and first order factors \(T - b_j\) with the \(b_j\) pairwise distinct elements of K. Deduce that there exists an element \(y_0\) in R such that \(f(y_0) = 0\) and such that \(x_0 - z\) and \(y_0 - z\) have the same sign, in E and R respectively, for all \(z \in \mathbb{K}\) (cf.~VI, p.~43, Ex. 26, c)).
\item
  Deduce from a) that there exists a K-isomorphism from the ordered field E onto a subfield of R. (First show that there exists such an isomorphism for the subfield \(\mathrm{K}(x_{0})\) of E; for any polynomial \(g \in \mathrm{K}[\mathrm{T}]\) such that \(\deg(g) < \deg(f)\) , apply the same argument used for \(f'\) in a) to evaluate the signs of \(g(x_{0})\) and \(g(y_{0})\) ; then apply Zorn's Lemma).
\item
  Show that if R and R' are two maximal ordered algebraic extensions of K then there exists precisely one K-isomorphism from R onto R' (apply b)).
\end{enumerate}

\begin{enumerate}
\def\labelenumi{\arabic{enumi})}
\setcounter{enumi}{15}
\tightlist
\item
  Let K be a maximal ordered field and let
\end{enumerate}

\[
g (\mathrm{T}) = a _ {0} + a _ {1} \mathrm{T} + \dots + a _ {n} \mathrm{T} ^ {n}
\]

be a nonzero polynomial in K{[}T{]} whose roots all lie in K. Show that for any polynomial \(f \in K[T]\) , the number of roots of the polynomial

\[
a _ {0} f (\mathrm{T}) + a _ {1} f ^ {\prime} (\mathrm{T}) + a _ {2} f ^ {\prime \prime} (\mathrm{T}) + \dots + a _ {n} f ^ {(n)} (\mathrm{T})
\]

which do not belong to K, counted with multiplicity, is at most equal to the number of roots of f, counted with multiplicity, which do not belong to K (use Ex. 12 for n = 1, then proceed by induction). Deduce that the polynomial

\[
a _ {0} + \frac {a _ {1}}{1 !} \mathrm{T} + \frac {a _ {2}}{2 !} \mathrm{T} ^ {2} + \dots + \frac {a _ {n}}{n !} \mathrm{T} ^ {n}
\]

has all its roots in K.

\begin{enumerate}
\def\labelenumi{\arabic{enumi})}
\setcounter{enumi}{16}
\tightlist
\item
  Let K be a maximal ordered field, let g be a nonzero polynomial in K{[}T{]} whose roots are all in K, and do not belong to the interval \([0, n]\) of K. Show that for every polynomial \(f(T) = a_{0} + a_{1}T + \cdots + a_{n}T^{n}\) of K{[}T{]}, the number of roots of the polynomial
\end{enumerate}

\[
a _ {0} g (0) + a _ {1} g (1) \mathrm{T} + a _ {2} g (2) \mathrm{T} ^ {2} + \dots a _ {n} g (n) \mathrm{T} ^ {n}
\]

which do not belong to K, counted with multiplicity, is at most equal to the number of roots of f, counted with multiplicity, which do not belong to K (same method as Ex. 16).

\begin{enumerate}
\def\labelenumi{\arabic{enumi})}
\setcounter{enumi}{17}
\item
  Let K be a maximal ordered field and let f be a polynomial of degree n in K{[}T{]}, all of whose roots are in K. Show that for all \(c \neq 0\) in K the polynomial \(f^{2} + cf'\) has at least n - 1 and at most \(n + 1\) roots in K, counted with multiplicity (use Ex. 13).
\item
  Let K be a maximal ordered field and \(f\) a nonzero polynomial in K{[}T{]}. A necessary and sufficient condition that, for every polynomial \(g \in \mathbf{K}[\mathbf{T}]\) , the number of roots of \(fg + g'\) , counted with multiplicity, not belonging to K, be at most equal to the number of roots of \(g\) , counted with multiplicity, not belonging to K, is that \(f(\mathbf{T}) = a - b\mathbf{T}\) with \(b \geqslant 0\) in K. (To see that the condition is sufficient use Ex. 12; to show that it is necessary, take g = 1 and g = f).
\item
  Let \(\mathbf{K}\) be a maximal ordered field and
\end{enumerate}

\[
f (\mathrm{T}) = a _ {0} + \binom {n} {1} a _ {1} \mathrm{T} + \binom {n} {2} a _ {2} \mathrm{T} ^ {2} + \dots + a _ {n} \mathrm{T} ^ {n}
\]

a polynomial in K{[}T{]}. For all integers p, q such that \(0 \leqslant p < p + q \leqslant n\) , the number of roots of the polynomial

\[
a _ {p} + \binom{q}{1} a _ {p + 1} \mathrm{T} + \binom{q}{2} a _ {p + 2} \mathrm{T} ^ {2} + \dots + a _ {p + q} \mathrm{T} ^ {q}
\]

counted with multiplicity, which do not belong to K is at most equal to the number of roots of f, counted with multiplicity, which do not belong to K (use Ex. 12).

Deduce that if \(b_{1}, \ldots, b_{n}\) are n distinct elements \textgreater0 in K, and if

\[
\left(\mathrm{T} + b _ {1}\right) \dots \left(\mathrm{T} + b _ {n}\right) = \mathrm{T} ^ {n} + \binom {n} {1} m _ {1} \mathrm{T} ^ {n - 1} + \dots + m _ {n}
\]

then

\[
m _ {k} ^ {1 / k} > m _ {k + 1} ^ {1 / (k + 1)} \quad \text { for } \quad 1 \leqslant k \leqslant n - 1.
\]

\begin{enumerate}
\def\labelenumi{\arabic{enumi})}
\setcounter{enumi}{20}
\tightlist
\item
  Let \((a_{i})_{1 \leqslant i \leqslant n}\) be a finite sequence of n elements of an ordered field K, not all zero; let \((a_{i_{k}})_{1 \leqslant k \leqslant p}\) be the subsequence of \((a_{i})\) consisting of the nonzero \(a_{i}\) ( \(i_{1} < i_{2} < \cdots < i_{p}\) ); then the number of indices \(k \leqslant p - 1\) such that \(a_{i_{k}}\) and \(a_{i_{k+1}}\) have opposite signs is called the number of variations of the sequence \((a_{i})\) .
\end{enumerate}

Let K be a maximal ordered field and let f be a nonzero polynomial in K{[}T{]} of degree n \textgreater{} 0; for all \(a \in K\) let \(w(a)\) denote the number of variations of the sequence \((f^{(i)}(a))_{0 \leq i \leq n}\) . If v is the number of roots of f, counted with multiplicity, in the interval \([a, b]\) (a \textless{} b), show that \(v \leq w(a) - w(b)\) and that \(w(a) - w(b) - v\) is even (« Budan-Fourier rule »; split the interval \([a, b]\) by the roots of f, and evaluate the amount of variation in \(w(x)\) as x passes one of these roots).

Deduce that if \(f(T) = a_{0} + a_{1}T + \cdots + a_{n}T^{n}\) , then the number of roots \textgreater{} 0 of f, counted with multiplicity, is at most equal to the number of variations of the sequence \((a_{i})_{0 \leq i \leq n}\) , and that the difference between these two numbers is even (« Descartes' rule »).

\begin{enumerate}
\def\labelenumi{\arabic{enumi})}
\setcounter{enumi}{21}
\tightlist
\item
  Let K be a maximal ordered field and let
\end{enumerate}

\[
f (\mathrm{T}) = a _ {0} + a _ {1} \mathrm{T} + \dots + a _ {n} \mathrm{T} ^ {n}
\]

be a polynomial in K{[}T{]} such that \(a_0 \neq 0 \neq a_n\) and

\[
a _ {p} = a _ {p + 1} = \dots = a _ {p + 2 m - 1} = 0 \quad (1 \leqslant p <   p + 2 m - 1 <   n).
\]

Show that \(f\) has at most \(n - 2m\) roots in \(\mathbf{K}\) , counted with multiplicity (apply Descartes' rule).

Deduce that if \(g(T) = 1 + c_{1}T + \cdots + c_{n}T^{n}\) has all its roots in K, and if \(h(T) = 1 + b_{1}T + \cdots + b_{2m}T^{2m}\) is the polynomial consisting of the first \(2m + 1\) terms of the formal power series \(1/g \in K[[T]]\) , then h has no roots in K.

\begin{enumerate}
\def\labelenumi{\arabic{enumi})}
\setcounter{enumi}{22}
\tightlist
\item
  Let K be a maximal ordered field and let \(a_{1}, \ldots, a_{n}, b_{1}, \ldots, b_{n}\) be 2n elements of K such that \(\sum a_{i} \neq 0\) and \(b_{1} < b_{2} < \cdots < b_{n}\) . Show that for any integer \(m \geqslant 1\) the number v of roots in K, counted with multiplicity, of the polynomial
\end{enumerate}

\[
f (\mathrm{T}) = a _ {1} \left(\mathrm{T} - b _ {1}\right) ^ {m} + a _ {2} \left(\mathrm{T} - b _ {2}\right) ^ {m} + \dots + a _ {n} \left(\mathrm{T} - b _ {n}\right) ^ {m}
\]

is at most equal to the number of variations w of the sequence

\[
\left(a _ {1}, a _ {2}, \dots , a _ {n}, (- 1) ^ {m} a _ {1}\right),
\]

and that the difference w - v is even. (Argue by induction on n, applying Ex. 13 (VI, p.~40) to a rational function of the form \(f(\mathrm{T})/(\mathrm{T}-c)^{m}\) .)

\begin{enumerate}
\def\labelenumi{\arabic{enumi})}
\setcounter{enumi}{23}
\item
  Let K be a maximal ordered field; consider an ordering on the field K(X) of rational functions which makes it an ordered extension of K. Show that this ordering is determined by the set A of \(x \in K\) such that x \textless{} X (show that the sign of every polynomial f over K is determined by A, using Prop. 9 of VI, p.~28). Conversely, show that to every set A ⊂ K such that \(y \in A\) whenever \(y \leqslant x\) and \(x \in A\) , there corresponds an ordering of K(X) making it an ordered extension of K and such that A is the set of \(x \in K\) such that x \textless{} X (same method). When A has a greatest element, or CA has a least element, or when A = K or A = ∅, then the corresponding ordering of K(X) is such that K(X) is not comparable to K. On the other hand, if neither A nor CA is empty, and there is no greatest element in A nor least element in CA, then K(X) is comparable to K.
\item
  Use Ex. 24 to find an example of a field E admitting several ordered field structures, none of which is induced from another via an automorphism of E. Hence find an example of a field, equipped with an ordering compatible with its ring structure but not a lattice ordering (consider the diagonal in E × E and use VI, p.~38, Ex. 5).
\item
  \begin{enumerate}
  \def\labelenumii{\alph{enumii})}
  \tightlist
  \item
    If K is an archimedean ordered field (VI, p.~40, Ex. 11, a)), and x and y are two elements of K with x \textless{} y, show that there exists \(r \in Q\) such that x \textless{} r \textless{} y. * Deduce that there is only one ordered subfield of R isomorphic to K. *
  \end{enumerate}
\end{enumerate}

\begin{enumerate}
\def\labelenumi{\alph{enumi})}
\setcounter{enumi}{1}
\tightlist
\item
  Consider the ordering of the field \(\mathbf{K} = \mathbf{Q}(\mathbf{X})\) in which \(\mathbf{X}\) is infinitely large with respect to \(\mathbf{Q}\) (Ex. 24). Show that \(\mathbf{K}\) is comparable to the subfield \(\mathbf{Q}(\mathbf{X}^2)\) , and give an example of two elements \(x, y\) of \(\mathbf{K}\) such that \(x < y\) and no element of \(\mathbf{Q}(\mathbf{X}^2)\) lies between \(x\) and \(y\) . c) Let \(\mathbf{K}\) be the ordered field defined in \(b\) ; show that the polynomial
\end{enumerate}

\[
f (\mathbf {Y}) = (\mathbf {Y} ^ {2} - \mathbf {X}) (\mathbf {Y} ^ {2} - 4 \mathbf {X}) - 1
\]

over K is irreducible, that it has roots in every maximal ordered extension of K, and that \(f(a) > 0\) for all \(a \in K\) .

\begin{enumerate}
\def\labelenumi{\arabic{enumi})}
\setcounter{enumi}{26}
\tightlist
\item
  Let K be an ordered field, let E be a purely transcendental extension of K and let \((\mathbf{X}_{i})_{i\in I}\) be a transcendence basis for E (V, p.~109).
\end{enumerate}

\begin{enumerate}
\def\labelenumi{\alph{enumi})}
\item
  * If K is archimedean then E has a structure of an ordered extension of K for which E is comparable to K if and only if the cardinality of I is at most equal to that of a transcendence basis for R over K (where K is considered as an ordered subfield of R, cf.~Ex. 26, a)); the set of such structures is then bijective with the set of injective maps from I into R such that \(f(I)\) is algebraically independent over K. *
\item
  If K is not archimedean then there always exists (at least) one structure on E of an ordered extension of K such that E is comparable to K (when I contains only one element, use Ex. 24 and notice that Q does not have a supremum in K; pass to the general case by Zorn's Lemma).
\end{enumerate}

* 28) a) Let K be a subfield of R and let \(\theta\) be a real number, algebraic over K. Show that the number of structures on K(θ) of an ordered extension of K is equal to the number of real conjugates of \(\theta\) (use Ex. 14, a) (VI, p.~40) and 26, a)).

\begin{enumerate}
\def\labelenumi{\alph{enumi})}
\setcounter{enumi}{1}
\tightlist
\item
  Let K be a subfield of R, equipped with n distinct structures of an ordered extension of Q, all archimedean; let \(P_{i}\) ( \(1 \leqslant i \leqslant n\) ) be the sets of positive elements under these orderings. Show that there exists a family of n elements \(b_{i}\) of K ( \(1 \leqslant i \leqslant n\) ) such that \(b_{i} \in P_{i}\) for all i and \(b_{i} \notin P_{k}\) for all \(k \neq i\) . (For each pair of distinct indices i, k, let \(a_{ik} \in P_{k}\) be such that \(a_{ik} \notin P_{i}\) ; consider the elements
\end{enumerate}

\[
c _ {i} = \sum_ {k \neq i} \left(\frac {1 + a _ {i k}}{1 - a _ {i k}}\right) ^ {2 r}
\]

where \(r\) is a sufficiently large integer.)

\begin{enumerate}
\def\labelenumi{\alph{enumi})}
\setcounter{enumi}{2}
\tightlist
\item
  Let \(\mathbf{K} = \mathbf{Q}(\theta)\) be an algebraic extension of \(\mathbf{Q}\) contained in \(\mathbf{R}\) , and let \(n\) be the number of real conjugates of \(\theta\) . Let \(\mathbf{C}\) be the set of sums of squares of elements of \(\mathbf{K}\) ; show that \(\mathbf{C}^* = \mathbf{C} \cap \mathbf{K}^*\) is a subgroup of \(\mathbf{K}^*\) and that \((\mathbf{K}^*: \mathbf{C}^*) = 2^n\) . (Use \(b\) ) and Cor. 2 of VI, p.~23.) *
\end{enumerate}

\begin{enumerate}
\def\labelenumi{\arabic{enumi})}
\setcounter{enumi}{28}
\item
  Let K be a maximal ordered field and let G be a subfield of K. Show that the set of extensions of G contained in K, and which are comparable to G, is inductively ordered; if \(E_{0}\) is a maximal element of this set, show that \(E_{0}\) is isomorphic to the field \(\mathrm{K}(\mathrm{G})\) defined in Ex. 11, b) of VI, p.~40 (prove that the natural map from \(\mathrm{F}(\mathrm{G})\) onto \(\mathrm{K}(\mathrm{G})\) maps \(E_{0}\) onto \(\mathrm{K}(\mathrm{G})\) , by first showing that \(E_{0}\) is a maximal ordered field, using Ex. 14, a) of VI, p.~40, and then that no element of \(\mathrm{K}(\mathrm{G})\) is transcendental over the image of \(E_{0}\) , using Ex. 24).
\item
  \begin{enumerate}
  \def\labelenumii{\alph{enumii})}
  \tightlist
  \item
    Let K be a maximal ordered field and let m and M be two elements of K with m \textless{} M. Show that every polynomial \(f \in K[X]\) which is positive in the interval \((m, M)\) is a sum of polynomials of the form \((\alpha X + \beta) g^{2}\) , where \(g \in K[X]\) and \(\alpha X + \beta\) is positive in \((m, M)\) . (Reduce to the case of first or second degree polynomials, and use the formulae
  \end{enumerate}
\end{enumerate}

\[
\begin{array}{l} (\mathbf {X} - a) (\mathbf {X} - b) = (\mathbf {X} - b) ^ {2} + (b - a) (\mathbf {X} - b) \\ (\mathbf {X} - a) (b - \mathbf {X}) = ((\mathbf {X} - a) (b - \mathbf {X}) ^ {2} + (b - \mathbf {X}) (\mathbf {X} - a) ^ {2}) / (b - a) \end{array}
\]

for \(a < b\) .)

\begin{enumerate}
\def\labelenumi{\alph{enumi})}
\setcounter{enumi}{1}
\tightlist
\item
  Show that the result of a) is no longer necessarily true when K is an arbitrary ordered field (notice that a polynomial can be \textgreater0 in K, but not in a maximal ordered extension of K; cf.~Ex. 26, c)).
\end{enumerate}

\begin{enumerate}
\def\labelenumi{\arabic{enumi})}
\setcounter{enumi}{30}
\tightlist
\item
  \begin{enumerate}
  \def\labelenumii{\alph{enumii})}
  \tightlist
  \item
    Let \(\mathbf{K}\) be a field whose algebraic closure \(\mathbf{E}\) is an extension of \(\mathbf{K}\) of prime degree \(q\) . Show that \(\mathbf{K}\) is perfect (V, p.~7).
  \end{enumerate}
\end{enumerate}

\begin{enumerate}
\def\labelenumi{\alph{enumi})}
\setcounter{enumi}{1}
\item
  Show that q is distinct from the characteristic of K (V, p.~162, Ex. 9).
\item
  Show that K contains the q-th roots of unity, and that E is the splitting field of an irreducible polynomial of the form \(X^{q}-a\) over K; deduce that q=2 (otherwise, use Ex. 5, V, p.~161 to deduce that \(X^{q^{2}}-a\) would be irreducible). Show also that -a is a square in K (V, loc. cit.), that -1 is not a square in K, and that \(E = K(i)\) ( \(i^{2} = -1\) ).
\item
  Now suppose that K is such that its algebraic closure E has arbitrary finite degree \(\neq1\) over K. Show that \(i\notin K\) and that \(\mathrm{E}=\mathrm{K}(i)\) (if \(\mathrm{E}\neq\mathrm{K}(i)\) , show by Galois theory that there would exist a field F such that \(\mathrm{K}(i)\subset\mathrm{F}\subset\mathrm{E}\) , and that E has prime degree over F; then apply c)). Deduce that K is a maximal ordered field (under some suitable ordering) (show by induction that any sum of n squares in K is a square in K; to prove that \(a^{2}+b^{2}\) is a square, consider a square root \(x+iy\) of \(a+ib\) in \(\mathrm{K}(i)\) ).
\end{enumerate}

\begin{enumerate}
\def\labelenumi{\arabic{enumi})}
\setcounter{enumi}{31}
\tightlist
\item
  Let A be an algebraically closed field of characteristic 0.
\end{enumerate}

\begin{enumerate}
\def\labelenumi{\alph{enumi})}
\item
  The only elements of finite order in the group \(\operatorname{Aut}(\mathbf{A})\) of automorphisms of \(\mathbf{A}\) are the identity element and elements of order 2 (involutions of \(\mathbf{A}\) ); these latter correspond bijectively to the maximal ordered subfields \(\mathbf{E}\) of \(\mathbf{A}\) such that \([\mathbf{A}:\mathbf{E}] = 2\) , that is \(\mathbf{A} = \mathbf{E}(i)\) (cf.~Ex. 31). If \(\sigma\) is an involution of \(\mathbf{A}\) and \(\mathbf{E}\) is the fixed subfield of \(\sigma\) , then the group \(\operatorname{Aut}(\mathbf{E})\) of automorphisms of \(\mathbf{E}\) is isomorphic to the quotient \(Z(\sigma) / \{e, \sigma\}\) , where \(Z(\sigma)\) is the centraliser of \(\sigma\) in \(\operatorname{Aut}(\mathbf{A})\) .
\item
  If A is algebraic over its prime field Q then any two involutions of A are conjugate in Aut(A) and the centraliser of an involution \(\sigma\) is \(\{e, \sigma\}\) ; the group Aut(A) and the set of involutions of A have the cardinality of the continuum, and the centre of Aut(A) is \(\{e\}\) .
\item
  If A has finite transcendence degree over Q, show that the centraliser \(Z(\sigma)\) of any involution \(\sigma\) of A is countable (observe that the fixed subfield E of \(\sigma\) is countable, and any automorphism of E which fixes the elements of some transcendence basis for E over Q is the identity). If deg. \(tr_{Q}A = n\) , show that there exist at least \(n + 1\) conjugacy classes of involutions of A (if \(E \subset A\) is maximal orderable with \([A : E] = 2\) , consider the transcendence degree over Q of the subfield of E consisting of elements fixed by Aut(E), and use Ex. 24).
\item
  If A has infinite transcendence degree m over Q, show that there exist involutions \(\sigma\) of A such that \(\operatorname{Card}(Z(\sigma)) = \operatorname{Card}(\operatorname{Aut}(A)) = 2^{m}\) (use Ex. 24 to show that there exist maximal orderable fields E such that deg. \(tr_{Q}E = m\) , such that for any subset M of any transcendence basis B for E over Q, there exists an automorphism \(u_{M}\) of E such that \(u_{M}(x) \neq x\) for \(x \in M\) and \(u_{M}(x) = x\) for \(x \in B - M\) ).
\item
  If \(\deg \cdot tr_{Q}A = 1\) , show that there is only one conjugacy class of involutions of A whose centralisers are infinite. (If E is a maximal orderable extension of Q of transcendence degree 1 such that \(\operatorname{Aut}(E)\) is nontrivial, use Ex. 24 to show that E cannot be comparable to Q; if \(t \in E\) is infinitely large with respect to Q, show that so is \(u(t)\) for any automorphism u of E, and that conversely if \(t'\) is infinitely large with respect to Q then there exists an automorphism u of E such that \(u(t) = t'\) , using Ex. 15 of VI, p.~40.)
\item
  Let u be an automorphism of A which commutes with all the involutions of A. Show that if \(a \in A\) is transcendental over Q then a and \(u(a)\) are algebraically dependent over Q. (Consider two elements b, c of A such that \(b^{2} = a\) , \(c^{2} = -u(a)\) and observe that there exists a maximal orderable extension \(E \subset A\) such that \([A : E] = 2\) and b, c are in E; obtain a contradiction by noticing that \(u(E) = E\) .)
\end{enumerate}

\(g\) ) Deduce from \(f\) ) that \(u(a) = a\) for every element \(a \in \mathbf{A}\) transcendental over \(\mathbf{Q}\) (consider a maximal ordered extension E of \(\mathbf{Q}(a)\) contained in A, algebraic over \(\mathbf{Q}(a)\) and comparable to \(\mathbf{Q}\) , using Ex. 24). Conclude that \(u\) is necessarily the identity (observe that if \(a\) is algebraic over \(\mathbf{Q}\) and \(b\) is transcendental over \(\mathbf{Q}\) then \(ab\) is transcendental over \(\mathbf{Q}\) ).

\begin{enumerate}
\def\labelenumi{\arabic{enumi})}
\setcounter{enumi}{32}
\tightlist
\item
  Let \(\mathbf{K}\) be a commutative field (of arbitrary characteristic \(q\) ), and let \(l\) be a prime number. Let \(\mathbf{K}'\) be an algebraic extension of \(\mathbf{K}\) with the following two properties:
\end{enumerate}

\begin{enumerate}
\def\labelenumi{(\roman{enumi})}
\item
  every polynomial in \(\mathbf{K}[\mathbf{X}]\) whose degree is not divisible by \(l\) has a root in \(\mathbf{K}'\) ;
\item
  every polynomial in \(\mathbf{K}'[\mathbf{X}]\) of degree equal to \(l\) has a root in \(\mathbf{K}'\) .
\end{enumerate}

\begin{enumerate}
\def\labelenumi{\alph{enumi})}
\item
  Show that the perfect closure \(\mathbf{K}_1\) of \(\mathbf{K}\) is contained in \(\mathbf{K}'\) and satisfies the condition analogous to (i).
\item
  Every algebraic extension of K whose degree is not divisible by l is isomorphic to a subextension of \(K'\) (use the primitive element theorem).
\item
  Let L be a subfield of \(K'\) containing K, and let M be a Galois extension of L, whose degree is a power of l. Show that M is isomorphic to a sub-extension of \(K'\) (use the fact that an l-group is nilpotent).
\end{enumerate}

\(d\) ) Deduce that \(\mathbf{K}'\) is an algebraic closure of \(\mathbf{K}\) (if \(\mathbf{F}\) is a (separable) extension of \(\mathbf{K}\) of finite degree, consider a Galois extension \(\mathrm{F}_1\supset \mathrm{F}\) of finite degree, and the fixed subfield of \(\mathrm{F}_1\) under a Sylow \(l\) -subgroup of the Galois group of \(\mathrm{F}_1\) over \(\mathbf{K}\) ).

\begin{enumerate}
\def\labelenumi{\arabic{enumi})}
\setcounter{enumi}{33}
\tightlist
\item
  Let E be a totally ordered set. A (Dedekind) cut is a partition (S, D) of E into two nonempty subsets such that D is of the form M(A), the set of upper bounds for some nonempty subset A of E; a necessary and sufficient condition for D to have a least element x is that \(D = \{x, \rightarrow\{, \text{so } S = \} \leftarrow, x\{.\)
\end{enumerate}

In a totally ordered abelian group G, a cut (S, D) is called proper if for all e \textgreater{} 0 in G there exist \(x' \in S\) and \(x'' \in D\) such that \(x'' - x' < e\) . Show that (S, D) is proper if and only if D is invertible in the monoid \(\mathfrak{M}(G)\) of major subsets of G (VI, p.~35, Ex. 30); all cuts are proper if and only if G is archimedean (VI, p.~35, Ex. 31); and hence * isomorphic to a subgroup of R (VI, p.~36, Ex. 33). *

\begin{enumerate}
\def\labelenumi{\arabic{enumi})}
\setcounter{enumi}{34}
\tightlist
\item
  A subfield K of an ordered field E is dense in E if for every nonempty open interval \(a, b \in E\) there exists \(x \in K\) such that a \textless{} x \textless{} b. Show that E is then comparable to K (Ex. 11, a)). If E is taken to be R(X), where R is a maximal ordered field and X is infinitely large with respect to R (VI, p.~43, Ex. 24), show that E is comparable to the subfield \(K = R(X^{2})\) , but that K is not dense in E.
\end{enumerate}

For K to be dense in E, it is necessary and sufficient that for all \(x \in E\) the partition of K consisting of \(K \cap J \leftarrow, x\) and \(K \cap [x, \rightarrow]\) be a proper cut (Ex. 34).

\begin{enumerate}
\def\labelenumi{\arabic{enumi})}
\setcounter{enumi}{35}
\item
  Let K be an ordered field and let \(\tilde{K}\) be the set of major sets D such that the cut (S, D) is proper; show that \(\tilde{K}\) is an ordered field, taking addition to be that of the monoid \(\mathfrak{M}(K)\) , and the product of two elements \(D_{1}\) and \(D_{2}\) of \(\tilde{K}\) contained in \(K_{+}\) to be the set \(\langle D_{1}D_{2}\rangle\) (VI, p.~35, Ex. 30), where \(D_{1}D_{2}\) denotes the set of \(x_{1}x_{2}\) for \(x_{1}\in D_{1}\) and \(x_{2}\in D_{2}\) . Show that for any increasing homomorphism f from K onto a dense subfield F of an ordered field E, there exists an increasing homomorphism g from E into \(\tilde{K}\) such that \(g\circ f\) is the natural injection of K into \(\tilde{K}\) .
\item
  Let K be an ordered field and let n \textgreater{} 0 be an integer such that every element x \textgreater{} 0 of K is equal to a power \(y^{n}\) of some element y \textgreater{} 0. Show that the field \(\tilde{K}\) has the same property.
\item
  Suppose that the field K is maximal ordered. Show that the field \(\tilde{\mathbf{K}}\) is also maximal ordered. (One can argue by contradiction, by considering an ordered algebraic extension E of \(\tilde{\mathbf{K}}\) which is maximal ordered, and supposing that \(\mathrm{E} \neq \tilde{\mathbf{K}}\) . Then there exists among the elements of \(\mathrm{E} \cap [\tilde{\mathbf{K}}\) an element \(w\) whose minimal polynomial \(f \in \tilde{\mathbf{K}}[\mathbf{X}]\) has least degree \(n > 1\) . Show that there exists an interval \(\mathbf{a}, b\) (in E such that \(a, b \in \mathbf{K}\) and \(a < w < b\) , and an element \(e > 0\) of E such that \(f'(x) > e\) for all \(x \in [\mathbf{a}, b]\) , or \(f'(x) < -e\) for all \(x \in [\mathbf{a}, b]\) . Show that there exists \(r > 0\) in K such that \(f(a) \leqslant -r\) and \(r \leqslant f(b)\) , or \(f(a) \geqslant r\) and \(f(b) \leqslant -r\) . Approximate the coefficients of \(f\) by elements of K, and show that we obtain in this way a polynomial \(g \in \mathbf{K}[\mathbf{X}]\) such that the function \(x \mapsto g(x)\) is monotone in \([\mathbf{a}, \mathbf{b}]\) and vanishes at some point of K in this interval, arbitrarily close to \(w\) .)
\item
  Let K be an ordered field and let E be an ordered algebraic extension which is maximal ordered. Show that the only K-automorphism of the field E is the identity. (First show that every K-homomorphism from E to itself is surjective, and is an isomorphism of ordered fields; if \(x \in E\) has minimal polynomial \(f \in K[X]\) then an automorphism of E necessarily permutes the roots of f belonging to E. Deduce that it fixes x.)
\item
  Let K be an ordered field. Make the field of formal power series \(E = K((X))\) an ordered field by taking the elements \textgreater0 to be the formal power series \(\sum_{n=-h}^{\infty} r_n X^n\) whose
\end{enumerate}

nonzero coefficient of least degree is \textgreater0.

\begin{enumerate}
\def\labelenumi{\alph{enumi})}
\item
  Show that if \(L = K(X)\) is the subfield of E consisting of rational functions in one indeterminate, then \(E = \tilde{L}\) (Ex. 36) up to isomorphism of ordered fields (observe that for every element a \textgreater{} 0 of L there exists an integer n \textgreater{} 0 such that \(0 < X^{n} < a\) ).
\item
  Suppose that K is maximal ordered. Show that the union F of the fields \(\mathbf{K}((\mathbf{X}^{1/n}))\) (ordered in the same way as E) is an ordered algebraic extension of E which is maximal ordered (use Ex. 2 of V, p.~150).
\item
  Show that the field \(\widetilde{\mathbf{F}}\) (Ex. 36) is the field of formal series \(\sum_{r} \alpha_r X^r\) , where \(r\) runs through
\end{enumerate}

the set Q of rational numbers, and where \(\alpha_{r} \in K\) and the set of r such that \(\alpha_{r} \neq 0\) is an increasing sequence which is either finite or tends to \(+\infty\) ; order this field in the same manner as E. In particular \(\widetilde{F} \neq F\) .

\begin{enumerate}
\def\labelenumi{\arabic{enumi})}
\setcounter{enumi}{40}
\tightlist
\item
  Let K be an ordered field, and E a vector space of dimension 2 over K.
\end{enumerate}

\begin{enumerate}
\def\labelenumi{\alph{enumi})}
\item
  Show that the set of triples of pairwise distinct lines ( \(D_{1}\) , \(D_{2}\) , \(D_{3}\) ) passing through 0 has two orbits under the action of the group \(GL^{+}(E)\) (consider the subgroup fixing two lines).
\item
  The set of triples of pairwise distinct half-lines ( \(\Delta_{1}\) , \(\Delta_{2}\) , \(\Delta_{3}\) ) with origin 0 has 7 orbits under the action of \(GL(E)\) , and 14 under the action of \(GL^{+}(E)\) (same method).
\end{enumerate}

\chapter{Modules over principal ideal domains}

\section{PRINCIPAL IDEAL DOMAINS}

\subsection{Definition of a principal ideal domain}

Recall (I, p.~104) that an ideal of a commutative ring A is said to be principal if it has the form \((a) = A a\) for some \(a \in A\) .

DEFINITION 1. --- A principal ideal domain is an integral domain (I, p.~116) in which every ideal is principal.

Examples. --- The ring Z of rational integers is a principal ideal domain (I, p.~111). If K is a commutative field, then the polynomial ring K{[}X{]} in one indeterminate over K is a principal ideal domain (IV, p.~11, Prop. 11); the same is true of the ring of formal power series K{[}{[}X{]}{]}, for every ideal of this ring has the form (X \(^{n}\) ) (IV, p.~38, Prop. 12). * The ring of integers in a p-adic field is a principal ideal domain. *

If \(\mathbf{Q}(i)\) denotes the field obtained from the field Q of rational numbers by adjoining a root i of the irreducible polynomial \(X^{2}+1\) , then the elements \(a+bi\) of \(\mathbf{Q}(i)\) , where a and b are rational integers, form a subring A of \(\mathbf{Q}(i)\) , called « the Gaussian integers », which is a principal ideal domain (VII, p.~50, Ex. 7). By contrast, in the field \(\mathbf{Q}(\rho)\) , where \(\rho\) is a root of \(X^{2}+5\) , the subring B consisting of elements \(a+b\rho\) (a and b rational integers) is not a principal ideal domain (VII, p.~51, Ex. 12).

The polynomial ring K{[}X, Y{]} in two indeterminates over a field K is not a principal ideal domain. Indeed the nonzero constants are the only elements which divide both X and Y, and none of them generates the ideal generated by X and Y.

\subsection{Divisibility in principal ideal domains}

Let A be a principal ideal domain and let K be its field of fractions (I, p.~116); we will see that the ordered group \(P^{*}\) of principal fractional ideals (VI, p.~6) of K is lattice ordered; more precisely:

PROPOSITION 1. --- Let K be the field of fractions of a principal ideal domain A, and let \((x_{\iota})_{\iota \in I}\) be a family of elements of K having a common denominator \(b \in K^{*}\) (in other words \(bx_{\iota} \in A\) for all \(\iota\) ). Then :

\begin{enumerate}
\def\labelenumi{\alph{enumi})}
\item
  The family \((x_{\iota})\) has a gcd in K.
\item
  Every gcd of \((x_{i})\) can be expressed in the form \(d = \sum_{i}a_{i}x_{i}\) , where the
\end{enumerate}

\(a_{i}\) are elements of A, all but finitely many of which are zero.

Indeed the ideal \(\sum_{i} A b x_{i}\) of \(A\) is principal, and so of the form \(A d'\) . Put

\(d' = bd (d \in \mathbf{K})\) . From the relation \(d' = \sum_{i} a_i b x_i\) , we deduce \(d = \sum_{i} a_i x_i\) , where

\(a_{i} \in A\) . Hence every common divisor of the \(x_{i}\) divides d.~On the other hand, since bd is a common divisor of the \(bx_{i}\) by construction, it follows that d is a common divisor of the \(x_{i}\) .

Remark. --- Prop. 1 applies with no restrictions to an arbitrary family \((x_{i})\) of elements of A (take \(b = 1\) ), and also to any finite family \((x_{i})\) of elements of K (if \(x_{i} = c_{i}b_{i}^{-1}\) with \(c_{i} \in \mathbf{A}\) and \(b_{i} \in \mathbf{A}\) , then take \(b\) to be the product of the \(b_{i}\) ).

COROLLARY. --- Let \((x_{t})\) be an arbitrary family of elements of a principal ideal domain A contained as a subring in an integral domain B, and let d be a gcd of the family \((x_{t})\) in A. Then the family \((x_{t})\) has gcd's in B, and d is one of them.

Indeed \(d\) is a common divisor of the \(x_{i}\) in B. On the other hand the relation \(d = \sum_{i} a_{i} x_{i}\) shows that every common divisor of the \(x_{i}\) in B divides \(d\) .

An important application of this corollary is when \(\mathbf{A} = \mathbf{K}[\mathbf{X}]\) and \(\mathbf{B} = \mathbf{E}[\mathbf{X}]\) , where \(\mathbf{K}\) is a field and \(\mathbf{E}\) is an extension of \(\mathbf{K}\) (IV, p.~12, Cor. 1).

The first assertion of Prop. 1 shows that the ordered group \(P^{*}\) is lattice ordered (VI, p.~10). In particular every finite family of elements of K admits an lcm. We can thus apply the results denoted (DIV) in VI, pp.~10 to 17 to principal ideal domains.

The following result is a consequence of the second assertion of Prop. 1 :

THEOREM 1 (« Bezout's identity »). --- Elements \(x_{\iota}\) ( \(\iota \in I\) ) of a principal ideal domain A are setwise coprime if and only if there exist elements \(a_{\iota}\) ( \(\iota \in I\) ) of A, all but finitely many of them zero, such that \(\sum a_{\iota}x_{\iota} = 1\) .

The necessity of the condition is just Prop. 1. Conversely, if \(\sum_{i} a_i x_i = 1\) then every common divisor of the \(x_i\) divides 1, and so 1 is a gcd of the \(x_i\) .

PROPOSITION 2. --- Let \(a, b, d, m\) and \(p\) be elements of the field of fractions \(K\) of a principal ideal domain \(A\) .

\begin{enumerate}
\def\labelenumi{\alph{enumi})}
\item
  « d is a gcd of a and b » is equivalent to « (d) = (a) + (b) ».
\item
  « m is an lcm of a and b » is equivalent to « m = (a) ∩ (b) ».
\item
  « p is an irreducible element of A » is equivalent to « (p) is a nonzero maximal ideal of A » and to « (p) is a nonzero prime ideal of A ».
\end{enumerate}

We have already proved \(a\) (Prop. 1). Since the common multiples of \(a\) and \(b\) are the elements of \((a) \cap (b)\) , and since, by hypothesis, the ideal \((a) \cap (b)\) is principal, say \((a) \cap (b) = (m)\) , it follows that \(m\) is an lcm of \(a\) and \(b\) , which proves \(b\) ). Finally, to say that \(p \neq 0\) is an irreducible element of \(A\) means by definition (VI, p.~17) that \((p)\) is a maximal element of the family of principal ideals \(\neq A\) of \(A\) , ordered by inclusion; since \(A\) has no ideals other than the principal ideals, this means that \((p)\) is a maximal ideal of \(A\) , whence \(c\) ), by the remark in VI, p.~17.

In a principal ideal domain A, the sum (resp. the intersection) of a finite number of ideals is also called the gcd (resp. lcm) of these ideals.

PROPOSITION 3. --- Let \(a, b, c\) be elements of the field of fractions of a principal ideal domain A, and let \(d\) be a gcd of \(a\) and \(c\) ; then the congruence \(ax \equiv b \pmod{c}\) has a solution \(x_0 \in A\) if and only if \(d\) divides \(b\) ; in this case the elements \(x\) of A which satisfy \(ax \equiv b \pmod{c}\) are precisely those which satisfy \(x \equiv x_0 \pmod{cd^{-1}}\) .

If \(ax \equiv b \pmod{c}\) , with \(x \in \mathbf{A}\) , then there exists \(y \in \mathbf{A}\) such that \(b = ax + cy\) , so \(d\) divides \(b\) . Conversely, if \(d\) divides \(b\) then \(b = ax_0 + cy_0\) with \(x_0, y_0 \in \mathbf{A}\) (Prop. 1), so \(ax_0 \equiv b \pmod{c}\) ; moreover, the relation \(ax \equiv b \pmod{c}\) is then equivalent to \(a(x - x_0) \equiv 0 \pmod{c}\) ; putting \(a = da'\) and \(c = dc'\) , the latter becomes \(a'(x - x_0) \equiv 0 \pmod{c'}\) . But this is equivalent (for \(x \in \mathbf{A}\) ) to \(x - x_0 \equiv 0 \pmod{c'}\) , since \(a'\) and \(c'\) are coprime (VI, p.~14, Prop. 10 (DIV) and VI, p.~15, Cor. 2 to Prop. 11 (DIV)).

PROPOSITION 4. --- Let \((a_i)_{1 \leq i \leq n}\) be a finite family of pairwise coprime elements of the principal ideal domain A. Then the canonical homomorphism (I, p.~110) from \(\mathbf{A}/\left(\prod_{i=1}^{n} a_i\right)\) into the product \(\prod_{i=1}^{n} \mathbf{A}/(a_i)\) is an isomorphism of A-algebras.

This follows from I, p.~110, Prop. 9 and from Prop. 2 a).

The conclusion of Prop. 4 is not valid if it is not assumed that the \(a_{i}\) are pairwise coprime (cf.~VII, p.~24, Prop. 9).

\subsection{Decomposition into irreducible factors in principal ideal domains}

We will now apply the results of VI, p.~18, relating to decomposition into irreducible elements, to principal ideal domains. By Prop. 2, an element \(p \neq 0\) of A is irreducible if and only if the ring \(\mathrm{A}/(p)\) is a field (I, p.~115, Cor. 1), that is if and only if the congruence \(ax \equiv b \pmod{p}\) admits a solution in A for each \(b \in \mathrm{A}\) and for each \(a \in \mathrm{A}\) which is not a multiple of \(p\) .

DEFINITION 2. --- Let A be an integral domain. A system of representatives of irreducible elements of A is a family \((p_{\alpha})\) of irreducible elements of A such that every irreducible element of A is an associate of precisely one \(p_{\alpha}\) .

THEOREM 2. --- Let A be a principal ideal domain and let \((p_{\alpha})\) be a system of representatives of irreducible elements of A. Then every nonzero element x of the field of fractions of A can be uniquely expressed in the form

\[
x = u \prod_ {\alpha} p _ {\alpha} ^ {n _ {\alpha}},\tag{1}
\]

where u is an invertible element of A, and where the \(n_{\alpha}\) are integers, all but finitely many of which are zero. For x to belong to A it is necessary and sufficient that all the \(n_{\alpha}\) be positive.

We will use the theorem about decomposition as a sum of irreducible elements (VI, p.~18, Th. 2), of which the above statement is only a translation. Since \(\mathcal{P}^*\) is a lattice ordered group it will be sufficient for us to show that every non empty set of principal ideals of A contains a maximal element, in order to check that the hypotheses of this theorem are indeed satisfied; now this follows from the next Lemma:

Lemma 1. --- Let A be a ring such that every left ideal of A is finitely generated. Then every nonempty set \(\Phi\) of left ideals of A, ordered by inclusion, has a maximal element.

By Zorn's Lemma (Set Theory, III, p.~154, Th. 2) it is enough to prove that \(\Phi\) is inductive. Now if \((\mathfrak{a}_{\lambda})\) is a totally ordered family of elements of \(\Phi\) then the union \(\mathfrak{a}\) of the ideals \(\mathfrak{a}_{\lambda}\) is a left ideal of \(\mathbf{A}\) , and so admits a finite system of generators \((a_i)_{1 \leqslant i \leqslant n}\) . Since each \(a_i\) belongs to an ideal \(\mathfrak{a}_{\lambda_i}\) , and since the family \((\mathfrak{a}_{\lambda})\) is totally ordered, the \(a_i\) all belong to the largest of the ideals \(\mathfrak{a}_{\lambda_i}\) , say \(\mathfrak{a}_{\mu}\) . Then \(\mathfrak{a} = \mathfrak{a}_{\mu}\) belongs to \(\Phi\) , which is thus indeed an inductive set.

Later we will study those rings B, called noetherian rings, such that every nonempty set of ideals of B contains a maximal element.

Remark. --- The family \((u, (n_{\alpha}))\) is called the decomposition of x into irreducible factors; by abuse of language, we also say that the formula (1) is the decomposition of x into irreducible factors. If \(x = u \prod_{\alpha} p_{\alpha}^{n_{\alpha}}\) and \(y = v \prod_{\alpha} p_{\alpha}^{m_{\alpha}}\) are

the decompositions of x and y into irreducible factors, then a necessary and sufficient condition for x to divide y is that \(n_{\alpha} \leqslant m_{\alpha}\) for all \(\alpha\) ; from this we deduce the formulae

\[
\operatorname * {g c d} (x, y) = \prod_ {\alpha} p _ {\alpha} ^ {\inf (n _ {\alpha}, m _ {\alpha})}\tag{2}
\]

\[
\operatorname{lcm} (x, y) = \prod_ {\alpha} p _ {\alpha} ^ {\sup \left(n _ {\alpha}, m _ {\alpha}\right)}.\tag{3}
\]

The property expressed by Th. 2 is true for a more general class of rings than principal ideal domains; we will study them later as unique factorisation domains; and we will see that polynomial rings and formal power series rings in arbitrarily many indeterminates are unique factorisation domains (Comm. Alg., VII, § 3).

\subsection{Divisibility of rational integers}

As was pointed out in section 1, the ring Z of rational integers is a principal ideal domain; its field of fractions is Q. The multiplicative group U of invertible elements of Z has two elements 1 and -1. The group \(Q_{+}^{*}\) of rational numbers \textgreater0 contains precisely one element from each class of associate elements of Q; it is thus isomorphic to the multiplicative group \(P^{*} = Q^{*}/U\) of principal fractional ideals of Q, with which it is usually identified. In particular, whenever gcd or lcm is used in the field Q (with respect to the ring Z), it is understood that these are elements \(\geqslant 0\) ; this convention allows us to speak of the gcd and the lcm of a family of rational numbers.

The irreducible integers \textgreater0 in Z are precisely those we have called prime numbers (I, p.~50) (they are sometimes called rational prime numbers); every irreducible element of Z is thus of the form p or -p, where p is a prime number, and the set P of prime numbers is a system of representatives of irreducible elements of Z.

PROPOSITION 5. --- The set of prime numbers is infinite.

Indeed, given an arbitrary finite family \((p_{i})\) \((1 \leqslant i \leqslant n)\) of distinct prime numbers, any prime divisor q of the number \(\left(\prod_{i=1}^{n} p_{i}\right) + 1\) (which is \textgreater1) is distinct from all the \(p_{i}\) , for otherwise it would divide 1.

\subsection{Divisibility of polynomials in one indeterminate over a field}

The polynomial ring K{[}X{]} in one indeterminate over a commutative field K is a principal ideal domain (IV, p.~11, Prop. 11). Its field of fractions is the field K(X) of rational functions in X with coefficients from K. The ring K{[}X{]} contains the subring of polynomials of degree 0, that is the field of constants, which is identified with K; the elements of K* are invertible in K, and hence in K{[}X{]}; conversely the formula \(\deg(uv) = \deg(u) + \deg(v)\) shows that every invertible polynomial has degree 0; the group U of invertible elements of K{[}X{]} is thus precisely K*. Thus two associate polynomials differ only by a nonzero constant factor; in particular every class of associate polynomials contains a unique monic polynomial. The subgroup of the multiplicative group K(X)* generated by the monic polynomials thus contains a unique element from each class of associate rational functions, and consequently is isomorphic to the group

\[
\mathcal {P} ^ {*} = \mathrm{K} (\mathrm{X}) ^ {*} / \mathrm{U}
\]

of principal fractional ideals of \(\mathrm{K}(\mathrm{X})\) . In particular, whenever gcd or lcm in the field \(\mathrm{K}(\mathrm{X})\) (with respect to the ring \(\mathrm{K}[\mathrm{X}]\) ) is mentioned, it will usually be understood that these are quotients of monic polynomials (or 0); this convention allows us to speak of the gcd or the lcm of a family of rational functions.

The irreducible elements of K{[}X{]} are precisely the irreducible polynomials in the usual sense (IV, p.~13, Def. 2), and the set of monic irreducible polynomials is a system of representatives of irreducible elements of K{[}X{]}.

A first degree polynomial is always irreducible. If K is an algebraically closed field then the converse is true (V, p.~19, Prop. 1); thus in this case every polynomial \(p(\mathbf{X})\) of degree \(n\) in \(\mathbf{K}[\mathbf{X}]\) can be written uniquely (up to the order of the factors) in the form

\[
p (\mathbf {X}) = c \left(\mathbf {X} - a _ {1}\right) \left(\mathbf {X} - a _ {2}\right) \dots \left(\mathbf {X} - a _ {n}\right)
\]

where \(c\) and the \(a_{i}\) are elements of \(\mathbf{K}\) .

PROPOSITION 6. --- For any field K, the set of monic irreducible polynomials in K{[}X{]} is infinite.

Indeed, given an arbitrary nonempty finite family \((p_{i})\) ( \(1 \leqslant i \leqslant n\) ) of distinct monic irreducible polynomials, the polynomial \(\left(\prod_{i=1}^{n} p_{i}\right) + 1\) is not invertible, and any monic irreducible factor q of this polynomial is necessarily distinct from all the \(p_{i}\) , otherwise it would divide 1.

\section{TORSION MODULES OVER A PRINCIPAL IDEAL DOMAIN}

\subsection{Modules over a product of rings}

Let \(\mathbf{A}\) be a ring and \((\mathfrak{b}_i)_{i\in I}\) a direct decomposition of \(\mathbf{A}\) , that is (I, p.~110, Def. 7) a finite family of two-sided ideals of \(\mathbf{A}\) such that the natural homomorphism from \(\mathbf{A}\) into the product of the \(\mathbf{A} / \mathfrak{b}_i\) is bijective. By loc. cit., Prop. 10, there exists a family \((e_i)_{i\in I}\) of central idempotents of \(\mathbf{A}\) such that \(\mathfrak{b}_i = \mathbf{A}(1 - e_i)\) , \(\sum_{i\in I}e_i = 1\) and \(e_i e_j = 0\) for \(i\neq j\) .

For every left A-module M, let \(M_{i}\) denote the set of \(m \in M\) such that \(b_{i}m = 0\) ; since \(b_{i}\) is a two-sided ideal, this is a submodule of M; moreover, if \(a, b \in A\) and \(a - b \in b_{i}\) , then the homotheties \(a_{M_{i}}\) and \(b_{M_{i}}\) coincide; there is thus a unique \((\mathbf{A}/\mathbf{b}_{i})\) -module structure on \(M_{i}\) such that the A-module structure of \(M_{i}\) is induced via the homomorphism \(A \to A/b_{i}\) .

PROPOSITION 1. --- The A-module M is the direct sum of its submodules \(M_{i}\) .

Let \(p_i: \mathbf{M} \to \mathbf{M}\) denote the homothety \(m \mapsto e_i m\) ; the map \(p_i\) is A-linear since \(e_i\) is central; since \(e_i^2 = e_i\) , \(\sum_{i \in I} e_i = 1\) and \(e_i e_j = 0\) for \(i \neq j\) , we have

\[
p _ {i} \circ p _ {i} = p _ {i}, \quad \sum_ {i \in I} p _ {i} = 1 _ {\mathrm{M}}, \quad p _ {i} \circ p _ {j} = 0 \quad \text { for } \quad i \neq j,
\]

and the \(p_i\) form an orthogonal family of projectors whose sum is the identity (II, p.~209, Def. 7). By loc. cit., Prop. 12, the module \(\mathbf{M}\) is the direct sum of the submodules \(p_i(\mathbf{M}) = e_i\mathbf{M}\) . In addition \(e_i\mathbf{M}\) is annihilated by \(\mathfrak{b}_i = \mathbf{A}(1 - e_i)\) ; if \(i \neq j\) and \(m \in \mathbf{M}\) then \((1 - e_i)e_jm = e_jm\) , whence no nonzero element of \(e_j\mathbf{M}\) is annihilated by \(1 - e_i\) , and so a fortiori by \(\mathfrak{b}_i\) . It follows that \(e_i\mathbf{M} = \mathbf{M}_i\) , and the proposition is proved.

Remarks. --- 1) Conversely, let \(\mathbf{M}_i'\) be an \((\mathbf{A} / \mathfrak{b}_i)\) -module for each \(i\) , and consider the A-module \(\mathbf{M}\) , the direct sum of the A-modules \(\mathbf{M}_i'\) ; then the submodules \(\mathbf{M}_i\) constructed above coincide with the \(\mathbf{M}_i'\) (it is enough to note that if \(i \neq j\) then no nonzero element of \(\mathbf{M}_j'\) is annihilated by \(\mathfrak{b}_i\) because \(\mathfrak{b}_i + \mathfrak{b}_j = \mathbf{A}\) ). Thus, roughly speaking, it amounts to the same thing to consider an A-module \(\mathbf{M}\) or a family \((\mathbf{M}_i)\) of modules over the rings \(\mathbf{A} / \mathfrak{b}_i = \mathbf{A}_i\) .

\begin{enumerate}
\def\labelenumi{\arabic{enumi})}
\setcounter{enumi}{1}
\item
  By the above proof, the projectors of M onto the components \(M_{i}\) are homotheties.
\item
  The A-module \(\mathbf{M}\) is cyclic if and only if each \(\mathbf{M}_i\) is cyclic: if \(\mathbf{M} = \mathbf{A}m\) , then \(\mathbf{M}_i = \mathbf{A}_i e_i m\) ; conversely if \(\mathbf{M}_i = \mathbf{A}_i m_i\) and \(m = \sum_{i \in I} m_i\) , then \(\mathbf{M} = \mathbf{A}m\) ; indeed, if
\end{enumerate}

\(n \in M\) projects onto \(a_i m_i\) for each \(i\) , and if \(a \in A\) is congruent to \(a_i\) mod \(b_i\) for each \(i\) , then \(am\) and \(n\) have the same image in each \(M_i\) , so coincide.

\begin{enumerate}
\def\labelenumi{\arabic{enumi})}
\setcounter{enumi}{3}
\tightlist
\item
  Let \(\mathbf{M}\) and \(\mathbf{N}\) be two A-modules, with components \((\mathbf{M}_i)\) and \((\mathbf{N}_i)\) . Let \(u \in \operatorname{Hom}_{\mathbf{A}}(\mathbf{M}, \mathbf{N})\) be an A-linear map from \(\mathbf{M}\) to \(\mathbf{N}\) ; then for all \(i\) and for all \(m \in \mathbf{M}_i\) we have \(u(m) \in \mathbf{N}_i\) , so \(u\) induces an \(\mathbf{A}_i\) -linear map \(u_i \in \operatorname{Hom}_{\mathbf{A}_i}(\mathbf{M}_i, \mathbf{N}_i)\) . It is easy to check that the map \(u \mapsto (u_i)\) is an isomorphism of \(\mathbf{Z}\) -modules (resp. of A-modules when \(\mathbf{A}\) is commutative)
\end{enumerate}

\[
\operatorname{Hom} _ {\mathbf {A}} (\mathbf {M}, \mathbf {N}) \rightarrow \prod_ {i \in I} \operatorname{Hom} _ {\mathbf {A} _ {i}} (\mathbf {M} _ {i}, \mathbf {N} _ {i}).
\]

\subsection{Canonical decomposition of a torsion module over a principal ideal domain}

Let M be a module over a commutative ring A. For each \(\alpha\in A\) let \(\mathbf{M}(\alpha)\) denote the kernel of the endomorphism \(x\mapsto\alpha x\) of M. If \(\alpha\) and \(\beta\) are two elements of A such that \(\alpha\) divides \(\beta\) , then clearly \(\mathbf{M}(\alpha)\subset\mathbf{M}(\beta)\) . In particular, as n runs through the set of rational integers \(\geqslant1\) , the submodules \(\mathbf{M}(\alpha^{n})\) form an increasing sequence; the union \(M_{\alpha}\) of the \(\mathbf{M}(\alpha^{n})\) is thus a submodule of M, consisting of those elements of M which are annihilated by some power of \(\alpha\) . For each submodule N of M, it is clear that \(N_{\alpha}=N\cap M_{\alpha}\) .

DEFINITION 1. --- Let \(\pi\) be an irreducible element of a principal ideal domain A; an A-module M is called \(\pi\) -primary if, for all \(x \in \mathbf{M}\) , there exists an integer \(n \geqslant 1\) such that \(\pi^n x = 0\) (in other words, if M is equal to the submodule \(\mathbf{M}_{\pi}\) ).

Clearly every cyclic module of the form \(\mathbf{A}/(\pi^{s})\) is \(\pi\) -primary. For an arbitrary A-module M, the submodule \(M_{\pi}\) is \(\pi\) -primary.

Lemma 1. --- Let M be a module over a principal ideal domain A; for all \(\alpha \in \mathbf{A}\) such that \(\alpha \neq 0\) , let \(\alpha = \varepsilon \sum_{i=1}^{r} \pi_i^{n(i)}\) be a decomposition of \(\alpha\) into irreducible factors (VII, p.~4). The submodule \(\mathbf{N} = \mathbf{M}(\alpha)\) of elements of M annihilated by \(\alpha\) is the direct sum of the submodules \(\mathbf{M}(\pi_i^{n(i)})\) , and the map which sends each \(x \in \mathbf{M}(\alpha)\) to its component in \(\mathbf{M}(\pi_i^{n(i)})\) has the form \(x \mapsto \gamma_i x\) ( \(\gamma_i \in \mathbf{A}\) ). Moreover

\[
\mathbf {M} \left(\pi_ {i} ^ {n (i)}\right) = \mathbf {N} \cap \mathbf {M} _ {\pi_ {i}} = \mathbf {N} _ {\pi_ {i}}.
\]

Note first that N is annihilated by \(\alpha\) , so has a natural \(\mathbf{A}/(\alpha)\) -module structure. By Prop. 4 of VII, p.~3, the canonical homomorphism from \(\mathbf{A}/(\alpha)\) into the product of the rings \(\mathbf{A}/(\pi_{i}^{n(i)})\) is a ring isomorphism; now applying Prop. 1 of VII, p.~6, we deduce that N is the direct sum of the \(\mathbf{M}(\pi_{i}^{n(i)})\) ; the projectors of this decomposition are homotheties, by VII, p.~7, Remark 2. The inclusion \(\mathbf{M}(\pi_{i}^{n(i)})\subset\mathbf{M}(\alpha)\cap\mathbf{M}_{\pi_{i}}\) is obvious; conversely, let \(x\in\mathbf{M}(\alpha)\cap\mathbf{M}_{\pi_{i}}\) . Then there exists a power \(\pi_{i}^{s}\) of \(\pi_{i}\) which annihilates x; we may assume \(s\geqslant n(i)\) ; by Bezout's identity, there exist \(\lambda,\mu\in\mathbf{A}\) such that \(\pi_{i}^{n(i)}=\lambda\pi_{i}^{s}+\mu\alpha\) , so \(\pi_{i}^{n(i)}x=0\) and so finally \(x\in\mathbf{M}(\pi_{i}^{n(i)})\) .

Lemma 2. --- Let M be a torsion module (II, p.~313) over an integral domain A. For every finite family \((x_{i})_{1 \leq i \leq n}\) of elements of M, there exists an element \(\gamma \neq 0\) in A such that the \(x_{i}\) all belong to \(\mathrm{M}(\gamma)\) .

Indeed, for each index \(i\) there exists an element \(\alpha_{i} \neq 0\) in A which annihilates \(x_{i}\) , and the element \(\gamma = \prod_{i=1}^{n} \alpha_{i}\) will fit the bill.

THEOREM 1. --- Let M be a torsion module over a principal ideal domain A; for each irreducible element \(\pi\) of A, let \(M_{\pi}\) be the submodule of M consisting of elements annihilated by some power of \(\pi\) . If P is a system of representatives of irreducible elements of A, then M is the direct sum of its submodules \(M_{\pi}\) for \(\pi \in P\) .

Every element \(x \in M\) belongs to the submodule \(\mathbf{M}(\alpha)\) for some \(\alpha \neq 0\) , so by Lemma 1 is a sum of a finite number of elements, each of which belongs to some submodule \(M_{\pi}\) . On the other hand, if \(\sum_{\pi \in P} x_{\pi} = \sum_{\pi \in P} y_{\pi}\) , where \(x_{\pi}, y_{\pi} \in M_{\pi}\) for all \(\pi \in P\) , and where all but finitely many of the \(x_{\pi}\) and \(y_{\pi}\) are zero, then Lemma 2 shows that there exists \(\gamma \neq 0\) in A such that all the \(x_{\pi}\) and \(y_{\pi}\) belong to the same submodule \(\mathbf{M}(\gamma)\) ; applying Lemma 1 to \(\mathbf{M}(\gamma)\) shows that \(x_{\pi} = y_{\pi}\) for all \(\pi \in P\) , which completes the proof.

Clearly, if \(\pi\) and \(\pi'\) are two associate irreducible elements, then \(M_{\pi} = M_{\pi'}\) ; thus, for a given module M, the submodule \(M_{\pi}\) depends only on the ideal ( \(\pi\) ) of A; it is called the \(\pi\) -primary component of the module M, and the decomposition of M as a direct sum of the \(M_{\pi}\) is called the canonical decomposition of M as a direct sum of its \(\pi\) -primary components.

COROLLARY 1. --- Every submodule N of a torsion module M is the direct sum of its submodules \(N \cap M_{\pi}\) .

This follows from the fact that \(\mathbf{N} \cap \mathbf{M}_{\pi}\) is the \(\pi\) -primary component \(\mathbf{N}_{\pi}\) of \(\mathbf{N}\) .

COROLLARY 2. --- The submodule N of the torsion A-module M is a direct factor if and only if \(N_{\pi}\) is a direct factor of \(M_{\pi}\) for every irreducible element \(\pi\) of A.

Indeed, if \(\mathbf{N}\) and \(\mathbf{N}'\) are two submodules of \(\mathbf{M}\) , then \(\mathbf{M} = \mathbf{N} \oplus \mathbf{N}'\) if and only if \(\mathbf{M}_{\pi} = \mathbf{N}_{\pi} \oplus \mathbf{N}_{\pi}'\) for every irreducible element \(\pi\) of \(\mathbf{A}\) (Cor. 1).

COROLLARY 3. --- Let N be a submodule of the torsion A-module M. If, for every irreducible element \(\pi\) of A, either \(N_{\pi} = 0\) or \((M/N)_{\pi} = 0\) , then N is a direct factor of M.

Indeed, the condition \((\mathbf{M}/\mathbf{N})_{\pi}=0\) implies \(N_{\pi}=M_{\pi}\) , and Cor. 2 applies.

An A-module M is called semi-simple if every submodule of M is a direct factor (cf.~A, VIII, § 3).

COROLLARY 4. --- Let A be a principal ideal domain which is not a field, and let M be an A-module. Then M is semi-simple if and only if M is a torsion module and \(\mathbf{M}_{\pi} = \mathbf{M}(\pi)\) for every irreducible element \(\pi\) of A.

First suppose that M is semi-simple; let \(x \in M\) and let \(\pi\) be an irreducible element of A. If N is a complement of \(A\pi x\) in M, then we can write \(x = \alpha\pi x + y\) , with \(\alpha \in A\) and \(y \in N\) ; but that implies \(y = (1 - \alpha\pi)x\) , so

\[
\pi (1 - \alpha \pi) x \in A \pi x \cap N = 0.
\]

It follows first of all that M is a torsion module; if moreover \(x \in M_{\pi}\) , then \(\pi(1 - \alpha\pi)x = 0\) , thus \(\pi x = \alpha\pi^{2}x = \alpha^{2}\pi^{3}x = \cdots = \alpha^{n}\pi^{n+1}x\) is zero and \(\mathbf{M}_{\pi} = \mathbf{M}(\pi)\) .

Conversely, by Cor. 2 it is enough to prove that an A-module M annihilated by an irreducible element \(\pi\) is semi-simple; but that is clear, since M then has a natural structure of a vector space over the field \(\mathrm{A}/(\pi)\) , and the submodules of M are precisely the vector subspaces under this structure.

Remark 1. --- Clearly the annihilator of every element \(\neq 0\) of a \(\pi\) -primary module has the form \(A\pi^k (k > 0\) an integer), since it is a principal ideal containing a power of \(\pi\) . Let \(x\) be an element of \(M\) ; for each \(\pi \in P\) , let \(x_{\pi}\) be the component of \(x\) in \(M_{\pi}\) ; the annihilator of \(x\) is the lcm of the annihilators of the nonzero \(x_{\pi}\) , but by the above it is equal in this case to the product of the annihilators of the nonzero \(x_{\pi}\) (VI, p.~16, Prop. 12 (DIV)).

PROPOSITION 2. --- If M is a finitely generated torsion module over a principal ideal domain A, then the \(\pi\) -primary components of M are zero except for a finite number of them, and the projectors of M onto these components \(M_{\pi}\) are homotheties.

This follows immediately from Lemma 1, for by Lemma 2 there exists \(\alpha \neq 0\) in A such that \(\mathbf{M} = \mathbf{M}(\alpha)\) .

Remark 2. --- By VII, p.~7, Remark 3, a finitely generated torsion A-module is cyclic if and only if each of its \(\pi\) -primary components is cyclic.

An important special case where Th. 1 and Prop. 2 apply is when A = Z; a Z-module is just an abelian group. An abelian group is called a torsion group if it is a torsion Z-module, that is if all its elements have finite order. Then P is taken to be the set of prime numbers \textgreater{} 0; for each prime number p \textgreater{} 0, an (abelian) group is said to be p-torsion if all its elements have orders which are powers of p.~In this terminology, Th. 1 shows that every torsion abelian group is a direct sum of p-torsion groups. In the case of finite groups, this also follows from I, p.~80, Th. 4.

\subsection{Applications : I. Canonical decompositions of rational numbers and of rational functions in one indeterminate}

THEOREM 2. --- Let A be a principal ideal domain, let K be its field of fractions and let P be a system of representatives of irreducible elements of A. Given an element \(x \in K\) , there exist a finite subset H of P, elements \(a_{0} \in A\) and \(a_{p} \in A\) not divisible by p in A ( \(p \in H\) ), and integers \(s(p) > 0\) ( \(p \in H\) ) such that

\[
x = a _ {0} + \sum_ {p \in \mathrm{H}} a _ {p} p ^ {- s (p)},\tag{1}
\]

where H and the \(s(p)\) are uniquely determined by these conditions.

Moreover, if \(R_{p}\) denotes a subset of A containing precisely one element of each residue class mod p in A ( \(p \in P\) ), then each \(x \in K\) can be uniquely expressed in the form

\[
x = a + \sum_ {p \in P} \left(\sum_ {h = 1} ^ {\infty} r _ {p h} p ^ {- h}\right)\tag{2}
\]

where \(a \in \mathbf{A}\) and \(r_{ph} \in \mathbf{R}_p\) for all \(h\) and \(p\) , and all but finitely many of the \(r_{ph}\) are zero.

Consider K as an A-module; then A is the submodule generated by 1. The quotient module K/A is the quotient of K by the equivalence relation \(x' - x \in \mathbf{A}\) , which is also written, in the notation of VI, p.~6, as \(x \equiv x' \pmod{1}\) ; let \(f\) denote the natural homomorphism from K onto M = K/A.

The quotient module M is a torsion module, for every element of M has the form \(f(a/b)\) ( \(a \in A\) , \(b \in A\) , \(b \neq 0\) ), whence \(bf(a/b) = f(a) = 0\) . Hence Th. 1 of VII, p.~8 applies. Let \(M_{p}\) ( \(p \in P\) ) be the submodule of elements of M annihilated by powers of p; then \(f^{-1}(M_{p})\) is the subring \(A_{p}\) of elements of K of the form \(ap^{-n}\) where \(a \in A\) and \(n \geqslant 0\) is an integer. The module M is the direct sum of the \(M_{p}\) , so every \(x \in K\) is congruent mod 1 to an element in the sum of the \(A_{p}\) ; in other words, formula (1) holds, with the \(s(p)\) integers \textgreater0, and the \(a_{p}\) finitely many elements of A such that \(a_{p}\) is not a multiple of p.

We now show that these conditions on the \(s(p)\) and the \(a_{p}\) completely determine H and the \(s(p)\) . Indeed H is then the set of \(p \in \mathbf{P}\) such that the component of \(f(x)\) in \(\mathbf{M}_{p}\) is \(\neq 0\) . On the other hand if \(s\) and \(s'\) are two integers \(>0\) such that \(s \geqslant s'\) and if \(a\) and \(a'\) are elements of \(A\) not divisible by \(p\) such that \(ap^{-s} \equiv a'p^{-s'} \pmod{1}\) , then we deduce that \(a \equiv a'p^{s - s'} \pmod{p^s}\) ; if \(s > s'\) then \(a \equiv 0 \pmod{p}\) , contradicting the hypotheses. This argument also shows that each \(a_p\) is well defined mod \(p^{s(p)}\) .

To complete the proof, we note first of all that in each residue class mod \(p^s\) in \(\mathbf{A}\) there exists a unique element of the form \(\sum_{h=0}^{s-1} r_h p^h\) with \(r_h \in \mathbb{R}_p\) for \(0 \leqslant h \leqslant s-1\) . In fact, we proceed by induction on \(s\) (the property follows from the definition of \(\mathbb{R}_p\) for \(s=1\) ): let \(x \in \mathbf{A}\) ; by hypothesis there is a unique element of the form \(\sum_{h=0}^{s-2} r_h p^h\) ( \(r_h \in \mathbb{R}_p\) ) in the residue class of \(x \mod p^{s-1}\) ; then \(x - \sum_{h=0}^{s-2} r_h p^h\) is a multiple \(ap^{s-1}\) of \(p^{s-1}\) ; now there exists a unique element \(r_{s-1}\) of \(\mathbb{R}_p\) such that \(a \equiv r_{s-1} (\bmod p)\) ; whence \(x \equiv \sum_{h=0}^{s-1} r_h p^h (\bmod p^s)\) . To obtain formula (2) it is now enough to apply this fact to each \(a_p\) in formula (1). The uniqueness is clear in view of the above.

The following are the most important applications of Th. 2:

I. The ring A is the ring Z of rational integers, and K = Q. Let P be the set of prime numbers \textgreater{} 0, and for each \(p \in P\) let \(R_{p}\) be the interval \((0, p - 1)\) in Z. Thus we have the canonical decomposition

\[
x = a + \sum_ {p \in P} \left(\sum_ {h = 1} ^ {\infty} e _ {p h} p ^ {- h}\right)
\]

with \(a \in \mathbf{Z}\) , \(e_{ph} \in \mathbf{Z}\) and \(0 \leqslant e_{ph} \leqslant p - 1\) .

\begin{enumerate}
\def\labelenumi{\Roman{enumi}.}
\setcounter{enumi}{1}
\tightlist
\item
  The ring A is the ring E{[}X{]} of polynomials in one indeterminate over a commutative field E, and \(\mathbf{K} = \mathbf{E}(\mathbf{X})\) . Let P be the set of monic irreducible polynomials in E{[}X{]} (VII, p.~5). For \(p \in P\) we can, by virtue of euclidean division of polynomials (IV, p.~10), take \(R_{p}\) to be the set of polynomials of degree strictly less than that of p.~Thus we have the decomposition (called canonical) of a rational function \(r(\mathbf{X}) \in \mathbf{E}(\mathbf{X})\) :
\end{enumerate}

\[
r (\mathrm{X}) = a (\mathrm{X}) + \sum_ {p \in \mathrm{P}} \left(\sum_ {h = 1} ^ {\infty} v _ {p h} (\mathrm{X}) \cdot p (\mathrm{X}) ^ {- h}\right)
\]

where \(a(\mathbf{X})\) is a polynomial and \(v_{ph}(\mathbf{X})\) is a polynomial of degree strictly less than that of \(p(\mathbf{X})\) , for all \(p\) and \(h\) . In particular, if the field \(\mathbf{E}\) is algebraically closed, then the \(p(\mathbf{X})\) have the form \(\mathbf{X} - \alpha\) with \(\alpha \in \mathbf{E}\) , and the \(v_{ph}(\mathbf{X})\) are thus constants.

Hence we can say that the vector space \(E(X)\) over the field E has a basis consisting of the monomials \(X^{n}\) ( \(n \geqslant 0\) an integer) and the rational functions of the form \(\mathbf{X}^{m}/(p(\mathbf{X}))^{h}\) , where \(p \in P\) and h and m are integers with \(h \geqslant 1\) and \(0 \leqslant m < \deg(p)\) .

\subsection{Applications : II. The multiplicative group of units of the integers modulo a}

Let \(a\) be a rational integer \(> 1\) , and let \((\mathbf{Z} / a\mathbf{Z})^*\) be the multiplicative group of invertible elements of the ring \(\mathbf{Z} / a\mathbf{Z}\) . If \(a = \prod_{i} p_i^{n(i)}\) is the decomposition of \(a\) into

prime factors, then the ring \(\mathbf{Z}/a\mathbf{Z}\) is isomorphic to the product of the rings \(\mathbf{Z}/p_{i}^{n(i)}\mathbf{Z}\) (VII, p.~3, Prop. 4) and the group \((\mathbf{Z}/a\mathbf{Z})^{*}\) is isomorphic to the product of the groups \((\mathbf{Z}/p_{i}^{n(i)}\mathbf{Z})^{*}\) . We are thus reduced to the study of the groups \((\mathbf{Z}/p^{n}\mathbf{Z})^{*}\) , where \(p\) is a prime number; recall (V, p.~80) that the order \(\varphi(p^{n})\) of \((\mathbf{Z}/p^{n}\mathbf{Z})^{*}\) is \(p^{n}-p^{n-1}=p^{n-1}(p-1)\) .

Suppose first of all that \(p > 2\) ; the natural homomorphism \(\mathbf{Z} / p^n\mathbf{Z} \to \mathbf{Z} / p\mathbf{Z}\) restricts to a homomorphism of groups from \((\mathbf{Z} / p^n\mathbf{Z})^*\) onto \((\mathbf{Z} / p\mathbf{Z})^*\) , whose kernel we denote \(\mathrm{U}(p^n)\) ; by VII, p.~3, Prop. 3 the residue class mod \(p^n\) of an integer \(m\) is invertible if and only if \(m\) is coprime to \(p\) , that is if and only if the residue class of \(m \mod p\) is invertible. It follows that \(\mathrm{U}(p^n)\) consists of all the residue classes mod \(p^n\) of integers congruent to \(1 \mod p\) , so has \(p^{n-1}\) elements, and that there is an exact sequence

\[
\{1 \} \rightarrow \mathrm{U} (p ^ {n}) \rightarrow (\mathbf {Z} / p ^ {n} \mathbf {Z}) ^ {*} \rightarrow (\mathbf {Z} / p \mathbf {Z}) ^ {*} \rightarrow \{1 \}.\tag{3}
\]

Similarly, for \(n \geqslant 2\) let \(U(2^n)\) denote the kernel of the natural homomorphism from \((\mathbf{Z}/2^n\mathbf{Z})^*\) to \((\mathbf{Z}/4\mathbf{Z})^*\) ; this is a group of order \(2^{n-2}\) , consisting of all the residue classes mod \(2^n\) of integers congruent to 1 mod 4, and there is an exact sequence

\[
\{1 \} \rightarrow \mathrm{U} (2 ^ {n}) \rightarrow (\mathbf {Z} / 2 ^ {n} \mathbf {Z}) ^ {*} \rightarrow (\mathbf {Z} / 4 \mathbf {Z}) ^ {*} \rightarrow \{1 \}.\tag{4}
\]

Lemma 3. --- Let \(x, y, k\) be integers with \(k \geqslant 0\) , and let \(p > 2\) be a prime number. If \(x \equiv 1 + py \mod p^2\) then \(x^{p^k} \equiv 1 + p^{k+1}y \mod p^{k+2}\) . If \(x \equiv 1 + 4y \mod 8\) then \(x^{2^k} \equiv 1 + 2^{k+2}y \mod 2^{k+3}\) .

To prove the first assertion, it is enough to show that, if \(k \geqslant 1\) and \(x \equiv 1 + p^{k}y \mod p^{k+1}\) , then \(x^{p} \equiv 1 + p^{k+1}y \mod p^{k+2}\) , and then to argue by induction on the integer k. For all \(a \in Z\) and \(k \geqslant 1\) , it is immediate that

\[
(1 + p ^ {k} a) ^ {p} \equiv 1 + p ^ {k + 1} a \mod p ^ {k + 2},
\]

hence

\[
\begin{array}{r l} (1 + p ^ {k} y + p ^ {k + 1} z) ^ {p} & = (1 + p ^ {k} (y + p z)) ^ {p} \equiv \\ & \equiv 1 + p ^ {k + 1} (y + p z) \equiv 1 + p ^ {k + 1} y \mod p ^ {k + 2} \end{array}
\]

Similarly, for \(k \geqslant 1\) we have

\[
(1 + 2 ^ {k + 1} a) ^ {2} \equiv 1 + 2 ^ {k + 2} a \mod 2 ^ {k + 3},
\]

so

\[
(1 + 2 ^ {k + 1} y + 2 ^ {k + 2} z) ^ {2} \equiv 1 + 2 ^ {k + 2} y \mod 2 ^ {k + 3},
\]

whence the second assertion by induction on k.

PROPOSITION 3. --- Let p \textgreater{} 2 be a prime number and let n \textgreater{} 0 be an integer ; then the group U(p\^{}n) is cyclic of order p\^{}\{n-1\} ; if n ≥ 2 then the residue class mod p\^{}n of an integer x congruent to 1 mod p is a generator of U(p\^{}n) if and only if x is not congruent to 1 mod p\^{}2. Let m \textgreater{} 1 be an integer ; then the group U(2\^{}m) is cyclic of order 2\^{}\{m-2\} ; if m ≥ 3 then the residue class mod 2\^{}m of an integer x congruent to 1 mod 4 is a generator of U(2\^{}m) if and only if x is not congruent to 1 mod 8.

Since \(\mathrm{U}(p^{n})\) has order \(p^{n-1}\) , the order of every element u of \(\mathrm{U}(p^{n})\) is a power of p, and u is a generator of \(\mathrm{U}(p^{n})\) if and only if \(u^{p^{n-2}} \neq 1\) . Now if u is the class of \(x = 1 + py\) , then \(u^{p^{n-2}}\) is the class of \(1 + p^{n-1}y\) , by Lemma 3, whence u generates \(\mathrm{U}(p^{n})\) if and only if \(y \neq 0 \mod p\) , in other words \(x \neq 1 \mod p^{2}\) . For example, the class \(1 + p\) generates \(\mathrm{U}(p^{n})\) . Similarly, the class u of x mod \(2^{n}\) generates \(\mathrm{U}(2^{n})\) if and only if \(u^{2^{n-3}} \neq 1\) , which means that x is not congruent to 1 mod 8, by Lemma 3; this is satisfied by x = 5.

Lemma 4. --- Let A be a principal ideal domain and let \(0 \rightarrow N \rightarrow M \rightarrow P \rightarrow 0\) be an exact sequence of A-modules. Suppose that there exist coprime elements a, \(b \in A\) such that aN = 0 and bP = 0. Then the exact sequence splits. If in addition N and P are both cyclic, then M is cyclic.

The module M is torsion, since abM = 0. The first assertion follows from Cor. 3 of VII, p.~9. If N and P are cyclic, then they are finitely generated, and hence so is M (II, p.~17, Cor. 5); since each p-primary component of M is isomorphic to a p-primary component either of N or of P, it follows from Remark 2 of VII, p.~10, that M is cyclic.

THEOREM 3. --- If \(a = \prod_{i} p_{i}^{n(i)}\) is the prime decomposition of the integer

a \textgreater{} 1, then the group \((\mathbf{Z}/a\mathbf{Z})^{*}\) of invertible elements of the ring Z/aZ is isomorphic to the product of the groups \((\mathbf{Z}/p_{i}^{n(i)}\mathbf{Z})^{*}\) . If p \textgreater{} 2 is a prime number and \(n \geqslant 1\) an integer, then the group \((\mathbf{Z}/p^{n}\mathbf{Z})^{*}\) is cyclic of order \(p^{n-1}(p - 1)\) . The group \((\mathbf{Z}/2\mathbf{Z})^{*}\) is trivial; for \(n \geqslant 2\) the group \((\mathbf{Z}/2^{n}\mathbf{Z})^{*}\) is the direct product of the cyclic group of order \(2^{n-2}\) generated by the residue class of 5 mod \(2^{n}\) and the cyclic group of order 2 consisting of the residue classes of 1 and -1 mod \(2^{n}\) .

The orders \(p^{n-1}\) of \(U(p^n)\) and \(p-1\) of \((\mathbf{Z}/p\mathbf{Z})^*\) are coprime; since \(U(p^n)\) and \((\mathbf{Z}/p\mathbf{Z})^*\) are cyclic (Prop. 3 and V, p.~78, Lemma 1), the group \((\mathbf{Z}/p^n\mathbf{Z})^*\) is cyclic (apply Lemma 4 to the exact sequence (3)). If \(n \geqslant 2\) then the restriction of the homomorphism \(v: (\mathbf{Z}/2^n\mathbf{Z})^* \to (\mathbf{Z}/4\mathbf{Z})^*\) to the subgroup \(\{1, -1\}\) is bijective; the group \((\mathbf{Z}/2^n\mathbf{Z})^*\) is thus the direct product of this subgroup and the kernel \(U(2^n)\) of \(v\) ; the result follows from Prop. 3.

Remark. --- Let p \textgreater{} 2 be a prime number and let x be an integer congruent to 1 mod p and not congruent to 1 mod \(p^{2}\) ; there is an exact sequence

\[
\{0 \} \rightarrow \mathbf {Z} / p ^ {n - 1} \mathbf {Z} \xrightarrow {u} (\mathbf {Z} / p ^ {n} \mathbf {Z}) ^ {*} \xrightarrow {v} (\mathbf {Z} / p \mathbf {Z}) ^ {*} \rightarrow \{1 \}\tag{5}
\]

of groups, where v is the natural projection and where u is induced on the quotients by the map \(r \mapsto x^{r}\) . Let \(Z_{p}\) be the ring of p-adic integers (V, p.~96), and let x be an element of \(Z_{p}\) such that \(x - 1 \in pZ_{p}\) and \(x - 1 \notin p^{2}Z_{p}\) ; on passing to inverse limits, the exact sequences (5) induce an exact sequence

\[
\{0 \} \rightarrow \mathbf {Z} _ {p} \stackrel {u} {\rightarrow} \mathbf {Z} _ {p} ^ {*} \stackrel {v} {\rightarrow} (\mathbf {Z} / p \mathbf {Z}) ^ {*} \rightarrow \{1 \}
\]

where \(v\) is the natural map, and the continuous map \(u\) extends the map \(n \mapsto x^n\) ( \(n \in \mathbf{Z}\) ). We often put \(x^n = u(x)\) for \(n \in \mathbf{Z}_p\) .

Similarly, if \(x \in \mathbf{Z}_2\) with \(x - 1 \in 4\mathbf{Z}_2\) and \(x - 1 \notin 8\mathbf{Z}_2\) , then there is a split exact sequence

\[
\{0 \} \rightarrow \mathbf {Z} _ {2} \stackrel {{u}} {{\rightarrow}} \mathbf {Z} _ {2} ^ {*} \stackrel {{v}} {{\rightarrow}} (\mathbf {Z} / 4 \mathbf {Z}) ^ {*} \rightarrow \{1 \},
\]

where \(u\) is a continuous extension of the map \(n \mapsto x^n\) .

\section{FREE MODULES OVER A PRINCIPAL IDEAL DOMAIN}

THEOREM 1. --- Let A be a ring such that every left ideal of A is projective (II, p.~231, Def. 1) as an A-module. Then every submodule M of a free left A-module L is a direct sum of modules isomorphic to ideals of A.

Let \((e_{\iota})_{\iota \in I}\) be a basis for L, and let \(p_{\iota}\) be the coordinate functions corresponding to this basis. Choose a well-ordering (Set Theory, III, p.~153) of I and let \(L_{\iota}\) denote the submodule generated by the \(e_{\lambda}\) for \(\lambda \leqslant \iota\) ; put \(M_{\iota} = M \cap L_{\iota}\) . The coordinate function \(p_{\iota}\) maps \(M_{\iota}\) onto an ideal \(\mathfrak{a}_{\iota}\) of A; since \(\mathfrak{a}_{\iota}\) is a projective A-module, there exists (II, p.~231, Prop. 4) a submodule \(N_{\iota}\) of \(M_{\iota}\) such that the map \(x \mapsto p_{\iota}(x)\) from \(N_{\iota}\) into \(\mathfrak{a}_{\iota}\) is bijective. Let \(M_{\iota}'\) be the submodule of L generated by the \(N_{\lambda}\) for \(\lambda \leqslant \iota\) ; we will show that \(M_{\iota}' = M_{\iota}\) for all \(\iota\) , which will imply that M is generated by the family \((N_{\iota})_{\iota \in I}\) . In fact, suppose \(M_{\lambda}' = M_{\lambda}\) for all \(\lambda < \iota\) ; then \(p_{\iota}(x) \in \mathfrak{a}_{\iota}\) for all \(x \in M_{\iota}\) ; thus there exists \(y \in N_{\iota}\) such that \(x - y\) is a linear combination of finitely many elements \(e_{\lambda}\) with \(\lambda < \iota\) ; in other words \(x - y\) is an element of \(M_{\lambda}\) for some \(\lambda < \iota\) ; the inductive hypothesis shows that \(x - y \in M_{\lambda}' \subset M_{\iota}'\) , that is that \(x \in M_{\iota}'\) , and so \(M_{\iota}' = M_{\iota}\) . It remains to show that the sum of the \(N_{\iota}\) is direct; now suppose there exists a linear relation \(\sum_{\iota} a_{\iota} = 0\) , with

\(a_{\iota} \in \mathbf{N}_{\iota}\) , where the \(a_{\iota}\) (all but finitely many of which are zero) are not all zero. Let \(\mu\) be the greatest index \(\iota\) such that \(a_{\iota} \neq 0\) ; since \(p_{\mu}(a_{\lambda}) = 0\) for \(\lambda < \mu\) , we have \(p_{\mu}(a_{\mu}) = p_{\mu}\left(\sum a_{\iota}\right) = 0\) , so \(a_{\mu} = 0\) , contradicting the choice of \(\mu\) .

COROLLARY 1. --- If every left ideal of A is projective, then every submodule of a projective left A-module is projective.

Indeed every projective A-module is a submodule of a free A-module (II, p.~231, Prop. 4), and Th. 1 applies.

COROLLARY 2. --- Every submodule of a free module over a principal ideal domain is free.

This follows immediately from Th. 1, since every ideal of a principal ideal domain is free.

COROLLARY 3. --- Every projective module over a principal ideal domain is free.

Remark. --- The proof of Th. 1 shows that every submodule of \(\mathbf{A}^{(1)}\) is isomorphic to a direct sum \(\oplus \mathfrak{a}_{\iota}\) , where each \(\mathfrak{a}_{\iota}\) is an ideal of \(\mathbf{A}\) .

PROPOSITION 1. --- If L is a free module of finite rank n over a principal ideal domain A, then every submodule M of L is a free module of rank ≤ n.

Indeed M is a free module by Cor. 2 to Th. 1, and it has rank \(\leqslant n\) by the previous remark, or by the following lemma:

Lemma 1. --- Let L be a module over a commutative ring A, generated by n elements, and let M be a free submodule of L; then M has rank ≤ n.

Suppose first that L is free. Let i denote the canonical injection of M into L. By III, p.~520, Cor., the homomorphism \(\bigwedge_{n+1}^{n+1} i: \bigwedge_{n+1}^{n+1} M \to \bigwedge_{n+1}^{n+1} L\) is injective; by III, p.~511, Prop. 6, \(\bigwedge_{n+1}^{n+1} L = \{0\}\) , so \(\bigwedge_{n+1}^{n+1} M = \{0\}\) ; it follows that M has rank \(\leqslant n\) (III, p.~518, Cor. 1). Now consider the general case; the module L is a quotient of a free module \(L'\) of rank n.~There exists a submodule \(M'\) of L isomorphic to M (II, p.~218, Prop. 21). By the first part of the argument \(M'\) has rank \(\leqslant n\) , and the result follows.

COROLLARY. --- Let E be a module over a principal ideal domain A, generated by n elements. Then any submodule F of E can be generated by at most n elements.

Indeed there exists a homomorphism \(f\) from \(\mathbf{A}^n\) onto \(\mathbf{E}\) (II, p.~216, Cor. 3), and \(f^{-1}(\mathbf{F})\) , which is a free module of rank \(m \leqslant n\) , is generated by \(n\) elements; the images of these elements under \(f\) generate \(\mathbf{F}\) .

\section{FINITELY GENERATED MODULES OVER A PRINCIPAL IDEAL DOMAIN}

\subsection{Direct sums of cyclic modules}

Let A be a commutative ring. Recall (II, p.~220, Prop. 22) that a cyclic A-module is isomorphic to a quotient module A/a, where a is an ideal of A. We will see later (Sect. 4) that every finitely generated module over a principal ideal domain is a direct sum of finitely many cyclic modules.

PROPOSITION 1. --- Let E be a module over a commutative ring A; suppose E is a direct sum of n cyclic modules \(A/\mathfrak{a}_{k}\) ( \(1 \leqslant k \leqslant n\) ), where the \(a_{k}\) are ideals of A;

then, for each integer \(p > 0\) , the A-module \(\bigwedge^p \mathbf{E}\) is isomorphic to the direct sum of the modules \(\mathbf{A} / \mathfrak{a}_{\mathrm{H}}\) , where for each \(p\) -element subset \(\mathbf{H} = \{k_1, \dots, k_p\}\) of \([1, n]\) , the ideal \(\mathfrak{a}_{\mathrm{H}}\) is \(\sum_{j=1}^{p} \mathfrak{a}_{k_j}\) .

Let \(x_{k}\) be the canonical generator of \(A / a_{k}\) , that is the image of the unit element of A, so that E is the direct sum of the \(Ax_{i}\) ( \(1 \leqslant i \leqslant n\) ). Then we know (III, p.~515, Prop. 10) that the exterior algebra \(\bigwedge_{n} E\) is isomorphic as an A-module to the tensor product \(\otimes_{i=1}^{n} (\bigwedge(Ax_{i}))\) . Now \(\bigwedge(Ax_{i})\) is just the direct sum \(A \oplus Ax_{i}\) , since

\[
\bigwedge^ {p}
\]

the exterior product of any two elements of \(Ax_{i}\) is zero, and so \(\Lambda E\) is the direct sum of the modules \(\mathbf{M}_{\mathrm{H}} = (\mathbf{A}x_{k_{1}}) \otimes \ldots \otimes (\mathbf{A}x_{k_{p}})\) as \(H = \{k_{1}, \ldots, k_{p}\}\) runs over the set of p-element subsets of {[}1, n{]} (with \(k_{1} < \ldots < k_{p}\) ); now \(M_{H}\) is known to be isomorphic to \(A/a_{H}\) (II, p.~257, Cor. 4), which completes the proof.

We now see that, in the notation of Prop. 1, if the ideals \(\mathfrak{a}_k\) form an increasing sequence, they are completely determined by the module E. More precisely:

PROPOSITION 2. --- Let A be a commutative ring, and let E be a direct sum of n cyclic modules \(A/a_{k}\) , where the \(a_{k}\) satisfy \(a_{1} \subset a_{2} \subset \ldots \subset a_{n}\) . Then, for \(1 \leqslant p \leqslant n\) , the ideal \(a_{p}\) is the annihilator of \(\bigwedge^{p} E\) ; if \(a_{n} \neq A\) then \(\bigwedge^{p} E \neq 0\) for \(1 \leqslant p \leqslant n\) and \(\bigwedge^{m} E = 0\) for m \textgreater{} n.

Indeed, in the notation of Prop. 1, we have \(\mathfrak{a}_{\mathrm{H}} = \mathfrak{a}_{s(\mathrm{H})}\) , where \(s(\mathrm{H})\) denotes the greatest element of the subset H. Since \(s(\mathrm{H}) \geqslant p\) for every \(p\) -element subset H, and since \(s(\mathrm{H}) = p\) for \(\mathrm{H} = \{1, \dots, p\}\) , it follows that \(\mathfrak{a}_p\) is the intersection of \(\mathfrak{a}_{\mathrm{H}}\) as H varies over the set of \(p\) -element subsets of \([1, n]\) ; the ideal \(\mathfrak{a}_p\) is thus indeed the annihilator of \(\bigwedge^p E\) , by Prop. 1.

COROLLARY. --- In the notation of Prop. 2, if \(\mathfrak{a}_n \neq \mathbf{A}\) , and if E is also isomorphic to the direct sum of \(m\) cyclic modules \(\mathbf{A} / \mathfrak{a}_j'\) with \(\mathfrak{a}_1' \subset \mathfrak{a}_2' \subset \ldots \subset \mathfrak{a}_m' \neq \mathbf{A}\) , then \(m = n\) and \(\mathfrak{a}_k = \mathfrak{a}_k'\) for \(1 \leqslant k \leqslant n\) (« uniqueness of the \(\mathfrak{a}_k\) »).
